\documentclass[twoside]{article}

\usepackage{aistats2027}
% If your paper is accepted, change the options for the package
% aistats2027 as follows:
%
% \usepackage[accepted]{aistats2027}
%
% This option will print headings for the title of your paper and
% headings for the authors names, plus a copyright note at the end of
% the first column of the first page.

% We also include a `preprint' option for non-anonymous preprints.
% Change the options for the package aistats2027 as follows:
%
%\usepackage[preprint]{aistats2027}
%
% This option will print headings for the title of your paper and
% headings for the authors names, but does not print the copyright and
% venue note at the end of the first column of the first page.

% Submissions have to follow US letter format.
% In case your local TeX installation defaults to A4, you will need to
% set papersize explicitly by activating the following three lines:
% \special{papersize = 8.5in, 11in}
% \setlength{\pdfpageheight}{11in}
% \setlength{\pdfpagewidth}{8.5in}

% If you use the natbib package, activate the following three lines:
\usepackage[round]{natbib}
\renewcommand{\bibname}{References}
\renewcommand{\bibsection}{\subsubsection*{\bibname}}

% If you use BibTeX in apalike style, activate the following line:
%\bibliographystyle{apalike}

\usepackage{algorithm} 
\usepackage{algorithmic}
\usepackage{amsmath, amssymb, graphicx}
\usepackage{amsthm}
\usepackage{placeins}
\graphicspath{{figures/paper/}}
% \usepackage[numbers,sort&compress]{natbib}
\newtheorem{theorem}{Theorem}
\newtheorem{lemma}{Lemma}
% \newtheorem{proof}{Proof}
\newtheorem{example}{Example}
\newtheorem{proposition}{Proposition}
\newtheorem{remark}{Remark}
\newtheorem{corollary}{Corollary}
\usepackage{xcolor}
\usepackage[hyphens]{url}
\usepackage{hyperref}
\hypersetup{
    breaklinks=true, % Allows links to break across lines
    colorlinks=false,
    urlcolor=black
}

\newcommand{\matt}[1]{\textcolor{purple}{[*** Matt: #1 ***]}}
\newcommand{\farzad}[1]{\textcolor{orange}{[*** Farzad: #1 ***]}}


\begin{document}

% If your paper is accepted and the title of your paper is very long,
% the style will print as headings an error message. Use the following
% command to supply a shorter title of your paper so that it can be
% used as headings.
%
%\runningtitle{I use this title instead because the last one was very long}

% If your paper is accepted and the number of authors is large, the
% style will print as headings an error message. Use the following
% command to supply a shorter version of the author names so that
% they can be used as headings (for example, use only the surnames)
%
%\runningauthor{Surname 1, Surname 2, Surname 3, ...., Surname n}

\twocolumn[

\aistatstitle{Sampling from Unknown Sources to Approximate a Known Target}

\aistatsauthor{ Author 1 \And Author 2 \And  Author 3 }

\aistatsaddress{ Institution 1 \And  Institution 2 \And Institution 3 } ]

\begin{abstract}
  This paper studies how to sequentially sample from multiple unknown source populations to match a known target population using as few samples as possible without rejection. This problem arises in polling, medical trials, and in scientific discovery. We give an algorithm for the simplest case of two sources and show that it matches a lower bound up to constant factors. We extend this algorithm to multiple sources and to when sources have different costs. Experiments based on the 2023 National Health Interview Survey show the benefit of this approach over baselines.
\end{abstract}


\section{Introduction}
% the problem: a gripping example
Imagine you are tasked with recruiting individuals to test a promising new medical treatment. You are asked to ensure that the overall population of individuals approximates the target population of the treatment. Further, to keep costs down, you must recruit as few individuals as possible, until your approximation error of the target is below a specific threshold. You have at your disposal different sources to select individuals for the test: online recruitment, telephone interviews, door-to-door questioning, and other outreach sources. While the population of each source is unknown, each is likely biased. For instance, online recruitment may over-represent younger subjects, whereas telephone interviews may over-represent older ones. How should you select the sources to best represent the target?

% why it matters
This question of representative, sample-efficient study design is at the heart of many real problems in medicine \citep{sanathanan1991randomized}, polling \citep{gelman2016high}, and scientific discovery \citep{gao2022sample}. Continuing the above example, in medical treatment trials of type 1 diabetes research has found large disparities in ethnic representation, including in key trials required for FDA approval of commercial devices \citep{akturk2021inequity}. This is a matter of life and death for underrepresented populations: even though only 9.3\% of type 1 diabetes cases in the United States affect Black populations \citep{akturk2021inequity}, their mortality rates are twice as high as other ethnic groups \citep{saydah2017disparities}. 

% a recommendation
One potential solution recommended by a recent review is to use adaptive trial
designs \citep{goshrani2026role}. This review argues that `research funders
and regulatory agencies should incentivize the design and conduct of adaptive
trials, with embedded community involvement, and particularly targeting
underrepresented populations'. Given a target population of individuals, how
can you adaptively sample individuals from your sources to match that target?

% % TODO: why it is hard
% However, if we have access to
% % TODO: briefly, what people have done in the past, and why our work is different
% In large part, 


% TODO: our approah
The difficulty is that we cannot choose which individual will arrive; we can only choose a sampling source, whose population may itself be unknown. The algorithm must therefore learn which sources are useful while continually reacting to the representation error in the sample collected so far.

Building on work in multi-armed bandits, we first study the one-dimensional setting and show that a simple feedback rule achieves the optimal sampling rate. We then extend the approach to multiple dimensions and, when sources have different costs, give an algorithm that can maintain this rate while identifying low-cost sources.

% TODO: re-title sections as number of dimensions instead of number of sources
% TODO: rename T to n


% We study this question as a sequential decision problem. The data collector has access to several sampling sources, but the population of each source is unknown in advance. At 


% The goal is to collect a set of samples that closely approximates a known target population. To do so, at each time step, the collector selects one source, and observes a randomly-sampled individual from that source.




% The objective is to make the
% empirical composition of the final dataset as close as possible to a known target
% distribution. Thus, the problem can be viewed as a form of adaptive sampling or
% multi-armed bandit problem in which the sources are sampling methods and the
% observations reveal their unknown demographic distributions. Unlike a standard bandit problem, however, the objective is
% not to maximize cumulative reward. Instead, the data collector must choose the
% sources sequentially so that the empirical demographic distribution of the final
% sample,





\section{Related Work}

\paragraph{Adaptive stratified Monte Carlo.}
A related line of work studies how to estimate an integral accurately when only
a limited number of function evaluations are available. The domain is divided
into smaller regions, called \emph{strata}. If the variance of the function is
larger in one region, more samples are needed there to estimate its contribution
accurately. Since these variances are generally unknown, 
\cite{carpentier2011finite} treat the strata as arms in a multi-armed bandit
problem and use an upper-confidence strategy to learn which regions have larger
variance and should receive more samples. This line of work was later extended
to settings in which the allocation and the partition itself are adapted
\citep{carpentier2012adaptive,carpentier2013toward,carpentier2015adaptive}.

Our objective is different. These methods adapt the sampling allocation to
reduce the uncertainty in estimating an unknown integral. In our setting, the
target \(p_G\) is already known. After every observation, we can measure the
realized deviation of the current sample from this target and choose the next
sampling method to correct that deviation. Thus, even if all arm distributions
were known, our feedback rule would still remain adaptive.

\paragraph{Post-hoc reweighting and importance sampling.}
A second related approach is to collect observations from distributions that do
not match the desired target and correct the mismatch afterward by assigning
different weights to the observations. This is the idea behind importance
sampling. When several sampling distributions are available, multiple
importance sampling combines observations from these different sources with
suitable weights \citep{elvira2019generalized}. There is also a large literature
in which the sampling distributions themselves are adapted over time
\citep{bugallo2017adaptive,elvira2021advances,lu2023daisee}.

This is closely related to post-hoc reweighting. Our method instead tries to correct the imbalance during data
collection. If some part of the target becomes overrepresented, the algorithm
changes the sampling method so that future observations are more likely to come
from underrepresented parts. The goal is therefore to make the unweighted
empirical sample itself representative, rather than correcting the sample
afterward through weights.
\paragraph{Online balancing and covariate-adaptive randomization.}
Another related line of work uses feedback to keep a cumulative imbalance
small. In online vector balancing, a random vector first arrives and the
algorithm then chooses a sign \(+1\) or \(-1\) for that vector in order to
keep the signed sum small
\citep{bansal2020online,bansal2021discrepancy}. A similar idea appears in
clinical trials. Pocock and Simon \citep{pocock1975sequential} observe the
characteristics of each new patient and then favor the treatment assignment
that reduces imbalance in these characteristics across treatment groups. In
both cases, however, the random observation is seen before the balancing
decision is made. In our setting, the order is reversed: we first choose a
sampling method \(A_t\), and only then observe the random outcome \(Y_t\)
generated by that method.
\paragraph{Herding.}
The closest connection to our representation objective is herding. Herding
starts from a known target and sequentially chooses observations so that their
empirical moments approach the target moments \citep{welling2009herding}. Kernel
herding extends this idea to continuous spaces and can achieve an
\(O(T^{-1})\) error for suitable classes of functions, compared with the usual
\(O(T^{-1/2})\) scale of independent random sampling
\citep{chen2010supersamples}. Herding has also been connected to
conditional-gradient optimization for minimizing moment discrepancy
\citep{bach2012equivalence}.

The main difference is that herding can choose the next observation directly.
In our setting, we can only choose the next sampling method \(A_t\); the
observation \(Y_t\) is then generated randomly from the distribution associated
with that method. In the finite setting, deterministic herding can be viewed as
a special case in which each arm returns a fixed observation with probability
one. Our setting therefore adds stochastic outcomes, and in the bandit version
the arm distributions must also be learned online.
\paragraph{Sampling costs.}
The cost problem we consider is different from the sampling-budget and
computational-cost questions studied in the related work above. \cite{carpentier2012adaptive} fix the total number of samples and optimize how these samples are
allocated across regions, but samples from different regions do not have
different costs. Importance-sampling methods may also account for computational
costs, such as the number of proposal evaluations required by different
weighting schemes. In our setting, however, each sampling method has its own
acquisition cost \(\ell_a\), and the goal is to find a low-cost combination of
methods that can produce the target \(p_G\).




\section{Warm up: one group}
\label{sec:two-source}
We first consider the simplest setting of a population defined by a single group. In general, the smallest number of sources needed to match a target population $p_G$ is two: one that has a larger proportion of the group and one with a smaller proportion. This warm-up isolates the main feedback mechanism. In Section~\ref{sec:multi}, we extend it to multiple demographic groups and sampling sources, and in Section~\ref{sec:cost} we incorporate  sampling costs.

Fix a single dimension, or group, of the population (e.g., individuals below age 65), and define
\[
Y_t =
\begin{cases}
1, & \text{if individual $t$ belongs to the group},\\
0, & \text{otherwise}.
\end{cases}
\]
Let \(p_G\) denote the proportion of this group in the target population.
After \(T\) respondents, define the empirical proportion of the designated
group by
\begin{equation}
\label{phat}
\widehat p_T
:=
\frac{1}{T}\sum_{t=1}^T Y_t.
\end{equation}
Our goal is to make \(\widehat p_T\) as close as possible to \(p_G\).

Imagine there are two unknown sources: online sampling, denoted by \(O\), and phone sampling, denoted by \(P\). The corresponding probabilities of sampling an individual from the designated group are \(p_O\) and \(p_P\). We make the  following regularity assumption about these groups
\[
p_P < p_G < p_O,
\]
i.e., online sampling over represents the designated group, whereas phone sampling under represents it. Without this assumption an adaptive method is not necessary, we can only get as close to $p_G$ as the closest source. Importantly, the numerical values of \(p_P\) and \(p_O\) need not be known.

Let \((\mathcal F_t)_{t\geq0}\) denote the history of the sampling process. Let \(A_t\) denote the action of choosing a source to sampling from, i.e., $A_t \in \{O,P\}$. This action is chosen using only past information, so
\(A_t\) is \(\mathcal F_{t-1}\)-measurable. Conditional on the chosen sampling method, we have
% \[
% Y_t\mid \mathcal F_{t-1},\,A_t=a
% \sim
% \operatorname{Bernoulli}(p_a).
% \]
% Equivalently,
\[
\mathbb E[Y_t\mid \mathcal F_{t-1}, A_t]
=
p_{A_t}.
\]

Consider the cumulative representation error
\begin{equation}
    S_t := \sum_{i=1}^t (Y_i-p_G).
\end{equation}
When \(S_t>0\), the current sample contains too many members of the designated
group, so the algorithm chooses the sampling method that is less likely to
recruit members of that group. When \(S_t\leq 0\), the current sample contains
too few members of the designated group, so it chooses the method that is more
likely to recruit members of that group.

\begin{algorithm}[t]
\caption{Sign-Based Adaptive Sampling}
\label{alg:sign-based}
\begin{algorithmic}[1]
\STATE Initialize \(S_0=0\).
\FOR{\(t=0,1,\ldots,T-1\)}
    \IF{\(S_t>0\)}
        \STATE Set \(A_{t+1}=P\) and sample by phone.
    \ELSE
        \STATE Set \(A_{t+1}=O\) and sample online.
    \ENDIF
    \STATE Observe \(Y_{t+1}\in\{0,1\}\).
    \STATE Update $S_{t+1}=S_t+Y_{t+1}-p_G$.
\ENDFOR
\end{algorithmic}
\end{algorithm}

The rule requires only the directional knowledge \(p_P<p_G<p_O\): it neither
estimates the two source distributions nor solves for an optimal sampling
mixture. Nevertheless, it keeps the cumulative representation error close to
zero. Together with Proposition~\ref{prop:one-dimensional-lower-bound}, the theorem below shows that
the sign-based rule is rate-optimal: no adaptive sampling algorithm can achieve
a uniformly better order than \(T^{-1}\).
\begin{theorem}[One group]
\label{thm:one-dimensional-guarantee}
Assume that
\[
p_P<p_G<p_O,
\]
and let Algorithm~1 generate the sequence of respondents. Then there exist
constants \(\lambda>0\) and \(C<\infty\), depending only on
\(p_P,p_G,\) and \(p_O\), such that for every \(T\geq1\) and every \(k\geq0\),
\[
\mathbb P\bigl(|S_T|\geq k\bigr)
\leq
Ce^{-\lambda k}.
\]
Consequently, for every \(\delta\in(0,1)\), with probability at least
\(1-\delta\),
\[
\left|\widehat p_T-p_G\right|
\leq
\frac{1}{\lambda T}
\log\left(\frac{C}{\delta}\right).
\]
\end{theorem}

The theorem says that the feedback rule prevents the representation error
\(S_T\) from growing with the sample size. In fact, the probability that
\(|S_T|\) is large decreases exponentially in its size.

The \(T^{-1}\) dependence in Theorem~\ref{thm:one-dimensional-guarantee}
is optimal in general. Even an algorithm that knows both sampling
distributions cannot achieve a uniformly smaller order of representation
error.

\begin{proposition}[Optimality of the \(T^{-1}\) rate]
\label{prop:one-dimensional-lower-bound}
Assume that $0<p_P<p_G<p_O<1$, and define $q=\min\{p_P,1-p_P,p_O,1-p_O\}>0$. 
Then, for any adaptive sampling algorithm and every \(T\geq1\),
\[
\mathbb P\left(
\left|\widehat p_T-p_G\right|
\geq \frac{1}{2T}
\right)
\geq q.
\]
In particular, no adaptive algorithm can improve the
\(T^{-1}\) rate uniformly over this class of problems.
\end{proposition}

\subsection{Significance of the Sign-Based Algorithm}

A useful way to understand the strength of Algorithm~\ref{alg:sign-based}
is to compare it with the usual scale of stochastic fluctuations. In a random
walk, or more generally in a sum of independent or martingale differences, the
cumulative fluctuation after \(T\) steps is typically of order \(\sqrt T\).
After dividing by \(T\), this leads naturally to an accuracy of order
\(T^{-1/2}\).

Algorithm~\ref{alg:sign-based} achieves something much stronger.
Theorem~\ref{thm:one-dimensional-guarantee} shows that $S_T=O_{\mathbb P}(1)$, 
and therefore $\widehat p_T-p_G=O_{\mathbb P}(T^{-1})$. 
% The difference is substantial: ignoring constants, \(0.01\) accuracy requires
% on the order of \(10^4\) samples at the usual \(T^{-1/2}\) rate, but only on
% the order of \(10^2\) samples at the \(T^{-1}\) rate.
The following decomposition lemma explains why this is surprising.

\begin{lemma}[Decomposition identity]
\label{lem:representation-decomposition}
Let \(\widehat p_T\) be defined as in \eqref{phat}, and let $y_T:=N_O(T)/T$ 
denote the fraction of respondents collected through online sampling. Then
\begin{equation*}
% \label{decomp. lemm.}
\widehat p_T-p_G
=
\underbrace{
(y_T-x^\star)(p_O-p_P)
}_{\text{allocation error}}
+
\underbrace{
\frac{1}{T}\sum_{t=1}^T
\bigl(Y_t-p_{A_t}\bigr)
}_{\text{sampling error}}
\end{equation*}where
\[
x^\star
:=
\frac{p_G-p_P}{p_O-p_P}
\]
is the proportion of online samples satisfying
\[
x^\star p_O+(1-x^\star)p_P=p_G.
\]
\end{lemma}

The second term in the decomposition of $\widehat p_T-p_G$ is ordinary sampling noise. The
martingale central limit theorem shows that this noise itself still fluctuates
on the \(\sqrt T\) scale, see Appendix~\ref{app:martingale-clt} for a proof of this. The allocation decisions of Algorithm~\ref{alg:sign-based} cancel it, leaving only a bounded
remainder. This is what allows the final representation error to improve from
the usual \(T^{-1/2}\) rate to \(T^{-1}\).

\subsection{Alternative methods}

% \paragraph{Cheap \& discard.} A natural alternative, commonly used in survey sampling, is to use a single cheap sampling method and then discard respondents from groups that become overrepresented. This can correct the composition of the retained sample, but it may require throwing away a large number of observations. Since every collected response has a cost, this can be expensive even when the cheapest sampling method is used. We next give a simple phone--online example showing that an adaptive mixture of sampling methods can achieve the desired representation at lower cost than collecting a biased sample and correcting it through rejection.

% \begin{example}
% Consider a single group and two sampling methods. Online sampling has $p_O=0.8,\ell_O=\$1$, while phone sampling has $p_P=0.2,\ell_P=\$3$. Let the target proportion be $p_G=0.4$. 
% Suppose we want a final sample of \(1000\) respondents. To obtain a sample of $1000$ respondents with \(p_G=0.4\),
% we need $400$ respondents from the group and $600$ not from the group. 
% If we use only online recruitment, only $20$\%
% of respondents are not in the group. Therefore, to obtain the
% required $600$ non-group respondents, we need $600/0.2=3000$ online respondents on average. Among these, about $2400$ are in the group and $600$ are not. We retain only $400$ of the respondents within the group and all $600$ respondents not in the group, and the total cost is \$$3000$. 
% By contrast, our adaptive method uses approximately \(333\) online respondents
% and \(667\) phone respondents. The resulting expected cost is approximately \$$2333$. Thus, using an
% appropriate mixture of the two methods is cheaper than collecting only from
% the cheapest method and correcting the sample by rejecting.
% \end{example}

\paragraph{Post-hoc reweighting.}
Another common approach is to keep the biased sample and assign larger
weights to underrepresented groups. This can be effective when the goal is
to estimate population quantities, but it is different from our goal as it does not change the composition
of the recruited sample itself. This distinction is important in settings
such as clinical trials, where having participants match the target proportions is itself part of the objective. 
% Uneven weights may also
% increase the variance of downstream estimates.


\paragraph{Explore first, then fix the allocation.}
Another natural strategy is to first estimate \(p_O\) and \(p_P\), compute the
corresponding oracle allocation \(x^\star\), and then sample according to this
fixed proportion. This can make the allocation error small, but the sampling
error in Lemma~\ref{lem:representation-decomposition} still fluctuates on the
order of \(\sqrt{T}\), due to the martingale central limit theorem (more details in Appendix~\ref{app:martingale-clt}). The resulting representation error therefore remains
of order \(T^{-1/2}\). Algorithm~\ref{alg:sign-based} instead continues to
react to the realized sample, which is what allows it to achieve the
\(T^{-1}\) rate.



\section{Multiple groups}
\label{sec:multi}
We now extend the feedback mechanism to multiple groups and
multiple sampling sources. Suppose there are \(K\) groups
and \(m\) available sampling sources. Each respondent is represented by
\[
Y_t\in\{0,1\}^K,
\qquad
\|Y_t\|_0\le \ell_0,
\]
where \(\|Y_t\|_0\) counts the number of nonzero coordinates. This is a
natural assumption for demographic data, since each respondent belongs
to only a limited number of groups, such as one racial group and one
age group. The target population distribution is \(p_G\in\mathbb R^K\). Source \(a\) has an unknown mean vector \(\mu_a\), with
\[
\mathbb E[Y_t\mid\mathcal F_{t-1},A_t=a]=\mu_a.
\]
We assume that \(p_G\) lies in the interior of the convex hull of the
source means. This is the multidimensional analogue of
\(p_P<p_G<p_O\). If \(p_G\) lies outside this convex hull, no mixture of
the available sources can produce the target distribution, so no sampling
algorithm can make the representation error converge to zero. The
interiority assumption further ensures that, whenever the sample deviates
from the target, there is a source whose expected contribution can
correct that deviation. 

For example, in our NHIS experiment we take \(K=3\), corresponding to
Hispanic, non-Hispanic Black, and other adults with diagnosed diabetes.
Each sampling source is a recruitment channel, such as online recruitment,
landline outreach, Medicaid, or regional clinics. A recruited individual is
represented by one of the three standard basis vectors, and \(p_G\) is the
corresponding ethnic composition of the target population.

Define the cumulative representation deficit by
\[
S_t:=\sum_{s=1}^t(Y_s-p_G).
\]
If the source means were known, the natural feedback rule would choose the source
whose expected contribution points most strongly against the current deficit:
\[
A_{t+1}
\in
\operatorname*{arg\,min}_{a\in[m]}
\left\langle S_t,\mu_a-p_G\right\rangle.
\]
Since the source means are unknown, the problem has a natural
multi-armed bandit structure. Unlike classical multi-armed bandits,
however, there is no single best source: the source that is most
useful depends on the current representation deficit.

We therefore use a UCB-type rule. Fix a confidence level
\(\delta\in(0,1)\). For each source \(a\), let
\[
N_a(t):=\sum_{s=1}^t \mathbf 1\{A_s=a\}
\]
denote the number of times source \(a\) has been sampled by time \(t\).
For \(n\ge1\), let \(\widehat\mu_{a,n}\) denote the empirical mean of the
first \(n\) observations collected from source \(a\). 
We assign source \(a\), after \(n\) observations, the confidence radius
\[
r_a(n)
:=
\sqrt{
\frac{K}{2n}
\log\left(
\frac{\pi^2 m K n^2}{3\delta}
\right)
},
\qquad n\ge1,
\]
and set \(r_a(0)=+\infty\). 
With probability at least \(1-\delta\), these confidence bounds hold
simultaneously for every source and every sample size; see the proof of
Theorem~\ref{thm:multidimensional-guarantee} for details.

Sources that have been sampled fewer times have larger confidence radii.
The algorithm subtracts this radius from the estimated directional
correction and chooses the source with the smallest resulting index. This
encourages exploration of poorly understood sources while continuing to
correct the current representation error.

We introduce a parameter \(\beta\) to control the strength of this
exploration bonus. Our theoretical guarantee uses \(\beta=1\), we examine smaller values empirically.

\begin{algorithm}[t]
\caption{Multi-Source Adaptive Sampling}
\label{alg:ucb-sampling}
\begin{algorithmic}[1]
\STATE Fix $\delta, \eta$, let \(S_0=0, N_a(0)=0\) for every source \(a\).
\FOR{\(t=0,1,\ldots,T-1\)}

    \IF{\(S_t=0\)}
        \STATE Choose \(A_{t+1}\) according to any fixed rule.
    \ELSE
        \STATE Set $u_t= S_t/\|S_t\|$.
        \STATE For every source with \(N_a(t)\ge1\), define
        \[
            I_a(t)
                =
                \left\langle
                u_t,\widehat\mu_{a,N_a(t)}-p_G
                \right\rangle
                -
                \beta\,r_a\!\left(N_a(t)\right).
                        \]
        For an unsampled source, set \(I_a(t)=-\infty\).
        \STATE Choose
        \[
            A_{t+1}
            \in
            \arg\min_{1\le a\le m} I_a(t).
        \]
    \ENDIF

    \STATE Sample from \(A_{t+1}\), observe
    \(Y_{t+1}\in\{0,1\}^K\).

    \STATE Update $S_{t+1} = S_t+Y_{t+1}-p_G$.

    \STATE Update \(N_{A_{t+1}}(t+1)\) and the empirical mean
    \(\widehat\mu_{A_{t+1},N_{A_{t+1}}(t+1)}\).

\ENDFOR
\end{algorithmic}
\end{algorithm}

\begin{theorem}[Multiple groups]
\label{thm:multidimensional-guarantee}

Assume that every observation \(Y_t\in\{0,1\}^K\) satisfies
\(\|Y_t\|_0\le \ell_0\), and suppose \(p_G\) lies in the
interior of the convex hull of the source means with margin \(c>0\).
For categorical observations \(Y_t\in\{e_1,\ldots,e_K\}\), the interior
and margin are understood relative to the affine hull of the probability
simplex.

Let Algorithm~\ref{alg:ucb-sampling} with \(\beta=1\)  and the confidence
radius defined above generate the sequence of respondents.

Then, for every \(T\ge1\) and every \(\eta\in(0,1)\), with probability at
least \(1-\delta-\eta\),
\[
\left\|\widehat p_T-p_G\right\|
=
O\!\left(
\frac{mK^{3/2}}{c^2T}
\log\frac{mK}{\delta c}
+
\frac{K}{cT}
\log\frac{K}{\eta c}
\right),
\]
where \(O(\cdot)\) hides only universal numerical constants.
\end{theorem}

\begin{proposition}[Optimality of the \(T^{-1}\) rate]
\label{prop:multidimensional-guarantee}

For every fixed \(K\geq 1\), there exists a collection of sampling
distributions and a target \(p_G\) satisfying the interiority assumption
such that, for every adaptive sampling algorithm and every \(T\geq 1\),
\[
\mathbb P\left(
\|\widehat p_T-p_G\|_2
\geq \frac{1}{2T}
\right)
\geq q
\]
for some constant \(q>0\) independent of \(T\).
% Thus the \(T^{-1}\) dependence in Theorem~\ref{thm:multidimensional-guarantee} cannot be improved
% uniformly over the class of admissible instances.
\end{proposition}

\section{Accounting for Costs}
\label{sec:cost}
In practice, different polling or data-collection methods have different
costs, so a natural question is how to assemble a representative sample as
cheaply as possible. We assign each sampling method \(a\) a cost \(\ell_a\).

A natural cost-aware modification of the adaptive rule is to choose at each
round the source giving the largest expected correction per unit cost, namely
\[
A_{t+1}
\in
\operatorname*{arg\,min}_{a\in[m]}
\frac{\langle S_t,\mu_a-p_G\rangle}{\ell_a}.
\]
This rule only asks which source looks cheapest for correcting the error
\emph{right now}. But our real goal is different: we want the entire final
sample to have the right composition at the smallest possible total cost.
A choice that looks cheap at one round may force us to use expensive sources
later. Our algorithm therefore first asks which combination of sampling
methods can produce \(p_G\) at the lowest average cost, and then runs the
adaptive sampling rule using that combination. The following example shows
that the difference can be large.

\begin{example}[The myopic rule can overpay]
\label{ex:myopic}
Let \(K=2\), \(p_G=0.5\), and consider three sources:
\(A\) with \(p_A=0.4\) and \(\ell_A=1\),
\(B\) with \(p_B=0.9\) and \(\ell_B=1\), and
\(C\) with \(p_C=0.1\) and \(\ell_C=3\). 
For \(S_t>0\), the myopic rule prefers \(C\) to \(A\), since
\[
\frac{p_C-p_G}{\ell_C}
=
\frac{-0.4}{3}
<
-0.1
=
\frac{p_A-p_G}{\ell_A}.
\]
For \(S_t<0\), it chooses \(B\). Thus, the myopic rule samples only
from \(B\) and \(C\). 
To obtain \(p_G=0.5\) from these two sources, they must be sampled in
equal proportions, since $(1/2)(0.9)+ (1/2)(0.1)=0.5$. The resulting average cost is therefore $2$ per recruit. In contrast, the pair \(\{A,B\}\) can also produce \(p_G=0.5\).
Since both sources cost \(1\), any mixture of them costs exactly \(1\)
per recruit. Thus the myopic rule pays twice as much, even though at
each round it chooses the source giving the largest correction per unit
cost.
\end{example}

A related approach is Lyapunov drift-plus-penalty~\citep{neely2010stochastic}.
Neely's method gives a tradeoff between sampling cost and the size of the
representation imbalance, and its basic form assumes that the source
distributions are known. In our setting, the source distributions are unknown
and must be learned while sampling. A natural learned-means version would choose
\[
A_{t+1}\in\arg\min_a
\left\{\langle S_t,\widehat\mu_a-p_G\rangle+V\ell_a\right\}.
\]
Its standard tradeoff is roughly \(\widetilde O(V/T)\) representation error
versus \(O(1/V)\) excess average cost. Thus improving one degrades the other.
By contrast, on a \(K\)-source support with positive margin, the mixture
representing \(p_G\) is unique; our support-selection theorem can therefore
retain the \(T^{-1}\) representation rate while also controlling the cost of
the selected support.

\subsection{The Algorithm}

For the cost results below, we restrict to the categorical setting
\[
Y_t\in\{e_1,\ldots,e_K\}.
\]
This restriction is mainly notational. If respondents are described by
several binary attributes, we can instead treat each possible combination of
attributes as a separate category. For example, Black \(<65\), Black \(65+\),
non-Black \(<65\), and non-Black \(65+\) can be treated as four categories,
so that every respondent belongs to exactly one category.

Thus the source means and \(p_G\) lie in the probability simplex, and the
relative-interiority margin is taken within its affine hull.
Our cost-aware procedure has two stages: we first select a low-cost
collection of sources, and then run Algorithm~\ref{alg:ucb-sampling}
restricted to the selected sources. 
We restrict the selected support to contain exactly \(K\) sources.
With more than \(K\) sources, there may be several different mixtures
of the selected sources whose mean is \(p_G\).
Algorithm~\ref{alg:ucb-sampling} controls representation error but does
not optimize over these different feasible mixtures, and could therefore
converge to an allocation whose cost is larger than the minimum-cost
representation of \(p_G\) on that support.

With \(K\) sources and positive margin, the representation of \(p_G\)
is unique. Thus, once the support is selected, the target allocation
and its cost are uniquely determined. 
For every \(K\)-source set \(I\), let \(c(I)\) denote the
relative-interiority margin of \(p_G\) in
\(\operatorname{conv}\{\mu_a:a\in I\}\), namely, the radius of the largest
Euclidean ball centered at \(p_G\), taken within the affine hull of the
probability simplex, that is contained in this convex hull.

Whenever \(c(I)>0\), let \(q(I)\in\Delta_K\) denote the unique mixture
satisfying
\[
p_G=\sum_{a\in I}q_a(I)\mu_a.
\]
Define
\[
B(I):=\sum_{a\in I}\ell_a q_a(I).
\]
\paragraph{Choosing the margin \(c_0\).}
Algorithm~\ref{alg:cost-support} contains a parameter \(c_0\) that determines
how far inside the convex hull the target must lie for a support to be
considered admissible. A larger \(c_0\) makes the problem statistically
easier, since the required exploration scales as \(c_0^{-2}\), but it may
exclude potentially cheaper supports. Conversely, decreasing \(c_0\) allows
the algorithm to consider a larger collection of supports, some of which may
have lower cost, but requires substantially more exploration and leads to
larger finite-sample constants. Thus \(c_0\) controls a tradeoff between
statistical robustness and the ability to search for cheaper sampling
strategies.

For an empirically admissible support \(I\), define
\[
\widehat B(I)
:=
\min_{q\in\Delta_K}
\left\{
\sum_{a\in I}\ell_a q_a:
\sum_{a\in I}q_a\widehat\mu_a=p_G
\right\},
\]
where \(\Delta_K\) is the \(K\)-dimensional probability simplex.
\begin{algorithm}[t]
\caption{Minimum-Cost Source Selection}
\label{alg:cost-support}
\begin{algorithmic}[1]

\STATE Fix a margin $c_0$, let $\widehat I^\dagger=\varnothing, \widehat B^\dagger=+\infty$.

\STATE Sample each source \(a\in[m]\) exactly \(n\) times and compute its
empirical mean \(\widehat\mu_a\).

\FOR{each \(K\)-source set \(I\subseteq[m]\)}

    \STATE Compute the empirical interior margin \(\widehat c(I)\).

    \IF{\(p_G\in\operatorname{conv}\{\widehat\mu_a:a\in I\}\) and
    \(\widehat c(I)\ge 3c_0/4\)}

        \STATE Compute $\widehat B(I)$.

        \IF{$\widehat B(I)<\widehat B^\dagger$}
            \STATE $\widehat B^\dagger=\widehat B(I)$,
            $\widehat I^\dagger=I$.
        \ENDIF

    \ENDIF

\ENDFOR

\IF{\(\widehat I^\dagger=\varnothing\)}
    \STATE Return \textsc{Fail}.
\ELSE
    \STATE Return \((\widehat I^\dagger,\widehat B^\dagger)\).
\ENDIF

\end{algorithmic}
\end{algorithm}

% Figure 1 is defined here so that it floats to the top of the first page of Section~\ref{sec:experiments}.
\begin{figure*}[t]
\centering
\includegraphics[width=\textwidth]{fig_main}
\vspace{-6mm}
\caption{(a) Two sources, share \(65+\) (2{,}000 runs). (b) Sixteen
sources, ethnic target (250 runs); ETC tuned to \(T=10^5\). (c) Black/other
target, offline \(\$199\): cost per recruit, exploration included (200
runs); dotted: the cheapest support \(I^\star\).}
\label{fig:main}
\end{figure*}

% \begin{algorithm}[t]
% \caption{Minimum-Cost Support Selection}
% \label{alg:cost-support}
% \begin{algorithmic}[1]
% \STATE Initialize $\widehat I^\dagger=\varnothing$ and
% $\widehat B^\dagger=+\infty$.
% \STATE Sample each source \(a\in[m]\) exactly \(n\) times and compute its
% empirical mean \(\widehat\mu_a\).

% \FOR{each \(K\)-source set \(I\subseteq[m]\)}
%     \STATE Compute the empirical interior margin \(\widehat c(I)\).
%     % \IF{\(\widehat c(I)\ge 3c_0/4\)}
%     %     \STATE Compute
%     %     \[
%     %     \widehat B(I)
%     %     :=
%     %     \min_{q}
%     %     \sum_{a\in I}\ell_a q_a
%     %     \]
%     %     subject to
%     %     \[
%     %     q_a\ge 0,\qquad
%     %     \sum_{a\in I}q_a=1,
%     %     \qquad
%     %     \sum_{a\in I}q_a\widehat\mu_a=p_G.
%     %     \]
%     \IF{\(p_G\in\operatorname{conv}\{\widehat\mu_a:a\in I\}\) and
%     \(\widehat c(I)\ge 3c_0/4\)}
%         \STATE Compute \(\widehat B := \min_q\{\sum_{a\in I}\ell_aq_a:\ q\ge0,\
%         \sum_{a\in I}q_a=1,\ \sum_{a\in I}q_a\widehat\mu_a=p_G\}\).
%         \IF{$\widehat B < \widehat B^\dagger$}
%             \STATE $\widehat B^\dagger = \widehat B$, $\widehat I^\dagger = I$
%         \ENDIF
%     \ENDIF
% \ENDFOR

% \IF{$\widehat I^\dagger=\varnothing$}
%     \STATE Return \textsc{Fail}.
% \ELSE
%     \STATE Return $(\widehat I^\dagger,\widehat B^\dagger)$.
% \ENDIF
% % \STATE Select $\widehat I \in \operatorname*{arg\,min}_{I \mid \widehat B(I) \in \mathcal{C}} \widehat B(I)$
% % \[
% % \widehat I
% % \in
% % \operatorname*{arg\,min}_{\substack{|I|=K\\
% % \widehat c(I)\ge 3c_0/4}}
% % \widehat B(I).
% % \]

% \STATE Return $(\widehat I^\dagger, \widehat B^\dagger)$.
% \end{algorithmic}
% \end{algorithm}

\begin{theorem}[Identifiability of min-cost sources]
\label{thm:cost-support}
Assume that there exists at least one \(K\)-source support \(I\subseteq[m]\)
with $c(I)\geq c_0$. 
Let
\[
\ell_{\max}:=\max_{a\in[m]}\ell_a.
\]
Suppose each source is sampled \(n\) times in
Algorithm~\ref{alg:cost-support}, where
\[
n
\geq
\frac{8K}{c_0^2}
\log\frac{2mK}{\delta}.
\]

Then, with probability at least \(1-\delta\), the support
\(\widehat I\) returned by Algorithm~\ref{alg:cost-support} satisfies $c(\widehat I)\geq \frac{c_0}{2}$, 
and
\[
B(\widehat I)
\leq
\min_{\substack{|I|=K\\ c(I)\geq c_0}}
B(I)
+
\frac{3\ell_{\max}}{c_0}
\sqrt{
\frac{K}{2n}
\log\frac{2mK}{\delta}
}.
\]
\end{theorem}
The benchmark in Theorem~\ref{thm:cost-support} is the cheapest support
with margin at least \(c_0\), not the cheapest feasible mixture, which may
have a smaller margin and be harder to identify. 
% The \(mn\) observations collected during the exploration phase need not be
% discarded and may also be retained in the final sample.
\begin{theorem}[Cost-aware sampling]
\label{thm:cost-aware}
Suppose there are \(m\) sampling methods and that there exists at least one
\(K\)-source support \(I\) such that \(c(I)\ge c_0>0\). Sample each source
\(n\) times, with
\[
n\ge
\frac{8K}{c_0^2}
\log\frac{2mK}{\delta},
\]
and use Algorithm~\ref{alg:cost-support} to select \(\widehat I\). After selecting \(\widehat I\), initialize Algorithm~\ref{alg:ucb-sampling} and run restricted to the sources in \(\widehat I\)
for \(T\) additional samples. Let \(\widehat p_T\) denote the empirical
composition of these \(T\) post-exploration samples. Then, with high
probability,
\[
\left\|\widehat p_T-p_G\right\|
=
O\!\left(
\frac{K^2}{c_0^2T}
\log\frac{K}{\delta c_0}
+
\frac{1}{c_0T}
\log\frac{1}{\eta c_0}
\right).
\]

Moreover,
\[
B(\widehat I)
\le
\min_{\substack{|I|=K\\ c(I)\ge c_0}}
B(I)
+
\frac{3\ell_{\max}}{c_0}
\sqrt{
\frac{K}{2n}
\log\frac{2mK}{\delta}
}.
\]

If
\[
\overline B_T
:=
\frac{1}{T}
\left(
\sum_{a=1}^m n\ell_a
+
\sum_{t=1}^T \ell_{A_t}
\right)
\]
denotes the total sampling cost, including the initial exploration phase,
amortized over the \(T\) post-exploration samples, then
\[
\overline B_T
\le
B(\widehat I)
+
\frac{mn\ell_{\max}}{T}
+
O_{\mathbb P}\!\left(
\frac{\ell_{\max}}{c_0}
\sqrt{\frac{K}{T}}
\right).
\]
\end{theorem}
\section{Experiments}
\label{sec:experiments}

The target population is the \(3{,}294\) adults with diagnosed diabetes in
the 2023 National Health Interview Survey (NHIS)~\citep{nhis2023}. We
construct \(m=16\) recruitment sources, e.g., online health-information
seekers, landline users, Medicaid patients and regional clinics;
\(\mu_a\) is the group composition of source \(a\) in NHIS, and recruits
are simulated as independent draws from \(\mu_a\), so \(T\) counts
simulated recruits (Appendix~\ref{app:exp-data}). Online recruits cost
\(\$72\) and offline recruits \(\$199\), the medians
of~\citet{brogger2020online}. The \emph{oracle closed loop} knows every
\(\mu_a\) and corrects \(S_t\) adaptively. The \emph{oracle open loop}
knows every \(\mu_a\) and computes in advance the cheapest mixture
\(x^\star\) with \(\sum_a x^\star_a\mu_a=p_G\); it keeps the share of each
source within \(O(1/T)\) of \(x^\star_a\), but does not react to the
realized recruits. \emph{Explore-then-commit} (ETC) samples uniformly for
\(10\%\) of the horizon, then follows the cheapest mixture matching \(p_G\)
under its estimates, with \(\delta=0.05\) (Appendix~\ref{app:exp-impl}).

\paragraph{One group: feedback gives the \(T^{-1}\) rate.}
For the share aged \(65+\) (\(40.0\%\) of online and \(78.0\%\) of landline
recruits, \(47.1\%\) of the target), Algorithm~\ref{alg:sign-based} achieves
the \(T^{-1}\) rate (Figure~\ref{fig:main}(a)): for \(T\ge10^3\), the
cumulative imbalance \(\sum_{t \geq T} |S_{t}|\) stays between \(1.4\) and
\(1.6\). The oracle open loop has the usual \(T^{-1/2}\) rate; its error at
\(T=10^5\) is \(84\) times larger. The difference is feedback: under
Algorithm~\ref{alg:sign-based}, the allocation error and the sampling noise
have correlation \(-0.9999\), so the adaptive allocation nearly cancels the
sampling noise (Appendix~\ref{app:exp-two}).

\paragraph{Multiple groups.}
The ethnic target is \{Hispanic: \(19.0\%\), non-Hispanic Black:
\(15.5\%\), other: \(65.4\%\)\} (\(K=3\), \(c=0.082\) over all sixteen
sources); the groups are mutually exclusive, so \(\ell_0=1\).
Algorithm~\ref{alg:ucb-sampling} learns the \(\mu_a\) and approaches the
oracle closed loop (Figure~\ref{fig:main}(b)): \(\|S_T\|\) is \(10.6\) at
\(T=1{,}202\) and \(2.5\) at \(T=10^5\), against \(2.0\) for the oracle. The
theoretical radius is conservative: \(90\%\) of runs reach error \(0.01\)
after \(3{,}631\) recruits with radius \(r_a\), and after \(1{,}202\) with
\(0.3\,r_a\). Some exploration is needed, however: with no radius, a source
that looks unhelpful after its first few samples may never be sampled
enough to correct the estimate, and at least \(10\%\) of runs stall
\(0.03\) or more from \(p_G\). Finer targets are harder to match. From
Black/other (\(K=2\), \(c=0.128\)) to ethnicity\(\times\)age (\(K=6\),
\(c=0.0014\)), \(c\|S_T\|\) at \(T=10^5\) stays in \([0.07,0.25]\) for the
oracle and in \([0.12,0.42]\) for Algorithm~\ref{alg:ucb-sampling},
suggesting that the observed imbalance grows roughly as \(1/c\). After
\(785\) recruits, Algorithm~\ref{alg:ucb-sampling} is \(5.8\) times closer
to the target than uniform recruitment for ethnicity, but only \(2.7\)
times for ethnicity\(\times\)age (Appendix~\ref{app:exp-many}).

\paragraph{Cost-aware sampling.}
We use the Black/other target (\(15.5\%\) Black; \(K=2\), \(c=0.128\) over
all sources), on which the myopic rule falls into the trap of
Example~\ref{ex:myopic}. Online recruits are \(13.9\%\) Black and southern
clinics \(24.5\%\), so the mixture \(I^\star\) of the two matches the target
at \(\$91.27\) per recruit, the cheapest feasible mixture, with margin
\(0.023\). Because online recruits correct an excess of Black recruits only
slowly, the myopic rule, with estimated means, prefers Spanish-language
clinics (\(0.4\%\) Black) at \(\$199\): it recruits \(60\%\) from southern
and \(36\%\) from Spanish-language clinics and pays \(\$197\) per recruit at
\(T=10^5\). Figure~\ref{fig:main}(c) shows the cost per recruit of
Algorithm~\ref{alg:cost-support} (\(c_0=0.02\), \(n=1{,}600\) samples per
source) followed by Algorithm~\ref{alg:ucb-sampling}, which keeps the
exploration samples in the sample and uses them to estimate the means.
Exploration costs \(\$183\) per recruit (\(\$4.69\)M in total). Afterwards
the cost per recruit falls to \(\$123\) at \(T=10^5\) and \(\$106\) at
\(T=3\times10^5\), against \(\$197\)--\(\$198\) for the myopic rule, which
is cheaper only below \(T\approx10^4\), while its estimates are poor;
\(194\) of \(200\) runs select \(I^\star\) or \(\{\)portal, southern
clinics\(\}\) (\(\$95.88\)). The price is a larger error at the same \(T\),
since the cheaper support has a smaller margin (\(0.023\) against
\(0.128\)): at \(T=10^5\), the median error is \(3.5\times10^{-5}\) against
\(8\times10^{-6}\), and one run selects a support that misses \(p_G\). With
a budget of \(\$16\)M, the myopic rule is \(2.3\) times as accurate, while
Algorithm~\ref{alg:cost-support}\(\to\)\ref{alg:ucb-sampling} recruits
\(1.7\) times as many participants (Appendix~\ref{app:exp-cost}).

Overall, the experiment shows a tradeoff between representation accuracy,
exploration cost, and recruitment cost. Restricting recruitment to a cheaper
support almost halves the cost per recruit once exploration is amortized,
while the composition error remains of order \(T^{-1}\), but larger, because
the cheaper support has a smaller margin.

\paragraph{Limitations.}
Exploration in Algorithm~\ref{alg:cost-support} is expensive (\(\$4.69\)M
here) and pays off only for large samples, and a too short exploration
selects supports that cannot reach \(p_G\). Variables outside the target
need not match: after matching ethnicity, adults in fair or poor health are
overrepresented by \(4.1\) points (Appendix~\ref{app:exp-downstream}).
Variables that affect the endpoint belong in the target, at the price of a
smaller margin.

% %\paragraph{Setup.}
% We use adults with diagnosed diabetes in the 2023 National Health Interview
% Survey (NHIS) as our target population~\citep{nhis2023}, giving \(3{,}294\)
% respondents. We construct \(m=16\) recruitment sources. Each source represents a population that
% could be reached through a particular recruitment method, such as online
% health-information seekers, landline users, Medicaid patients, or patients
% at clinics in different regions (Appendix~\ref{app:exp-data}). For each source
% \(a\), \(\mu_a\) is the vector of group proportions in that source. For example, for the
% ethnicity target, \(\mu_a\) gives the proportions of Hispanic,
% non-Hispanic Black and other demographics within source \(a\).

% We assign a cost of \(\$72\) per recruit to the two online sources and
% \(\$199\) to the remaining sources, based on the online and offline recruitment
% costs reported in~\citep{brogger2020online}.

% \paragraph{Baselines.} 
% We compare our algorithms with several baselines. The \emph{oracle closed
% loop} knows every \(\mu_a\), and uses this to adaptively choose sources to correct the current representation error. The
% \emph{oracle open loop} also knows every \(\mu_a\), but instead computes in
% advance the cheapest fixed mixture \(x^\star\) satisfying
% \[
% \sum_a x^\star_a\mu_a=p_G.
% \]
% It then chooses sources so that, after \(T\) rounds, the fraction of samples
% drawn from each source is within \(O(1/T)\) of \(x^\star_a\). Importantly, this
% allocation is fixed in advance: if the realized recruits happen to
% overrepresent or underrepresent a group, the oracle does not change its
% mixture in response. 
% Finally, \emph{explore-then-commit} (ETC) does not know the source
% distributions. It samples all sources uniformly for the first \(10\%\) of the
% horizon, estimates their means, and then follows the cheapest mixture that
% matches \(p_G\) under those estimates. We use \(\delta=0.05\);
% additional details are given in Appendix~\ref{app:experiments}.

% \paragraph{One group: feedback gives the \(T^{-1}\) rate.}
% We first test the two-source setting of Section~\ref{sec:two-source} on
% NHIS data. The target has one group: adults aged \(65+\). We use two recruitment sources: \(40.0\%\) of online recruits
% are aged \(65+\), compared with \(78.0\%\) of landline recruits, while the
% target population is \(47.1\%\) aged \(65+\).

% Algorithm~\ref{alg:sign-based} achieves the predicted \(T^{-1}\) rate:
% for \(T\ge10^3\), the cumulative imbalance \(\sum_{t \geq T} |S_{t}|\) stays between \(1.4\)
% and \(1.6\). In contrast, the oracle open loop knows the true composition
% of both sources but follows a fixed mixture, and therefore achieves only the
% usual \(T^{-1/2}\) rate. At \(T=10^5\), its error is \(84\) times larger.

% The difference comes from feedback: Algorithm~\ref{alg:sign-based}
% continually adjusts the next source in response to the realized imbalance.
% Its allocation error and sampling noise have
% correlation \(-0.9999\), showing that the adaptive allocation nearly
% cancels the sampling noise. Further robustness experiments are reported in
% Appendix~\ref{app:exp-two}.

% \paragraph{Multiple groups.}
% We next consider a multi-group \(K\!=\!3\) race/ethnicity target corresponding to the proportions of individuals diagnosed with diabetes \{Hispanic: 19.0\%, non-Hispanic Black: 15.5\%, other: 65.4\%\}. We use all sixteen possible recruitment sources. 
% % \{Hispanic,
% % non-Hispanic Black, other\}, using all sixteen recruitment sources.
% Since these three groups are mutually exclusive, each recruit belongs to
% exactly one group, so \(\ell_0=1\).
% The target has margin \(c=0.082\)
% (Figure~\ref{fig:main}(b), Appendix Table~\ref{tab:many}).

% The main question is whether Algorithm~\ref{alg:ucb-sampling}, which must
% learn the unknown source distributions, can approach the performance of the
% closed-loop oracle that knows them from the start. At \(T=1{,}202\), the
% algorithm is still learning and has cumulative imbalance
% \(\|S_T\|=10.6\). By \(T=10^5\), this falls to \(2.5\), compared with
% \(2.0\) for the oracle. Thus, after the initial learning period, the
% unknown-means algorithm performs nearly as well as the known-means oracle.

% Exploration is important during this learning period. If we set
% \(r_a\equiv0\) (no exploration bonus at all), the algorithm trusts its
% current estimates completely. A source that is poorly estimated from its
% first few samples may therefore be used too little afterward, and some runs
% become stuck far from the target. Indeed, at least \(10\%\) of runs finish
% \(0.03\) or more from \(p_G\).

% The theoretical confidence radius is conservative. With the full radius
% \(r_a\), \(90\%\) of runs reach error at most \(0.01\) after \(3{,}631\)
% recruits. Using the smaller radius \(0.3\,r_a\), the same criterion is
% reached after only \(1{,}202\) recruits. Thus, in this experiment, reducing
% the weight placed on exploration leads to substantially faster convergence.
% Some exploration is necessary, but the full theoretical exploration bonus
% appears to be larger than needed in practice.

% \paragraph{The margin decides what can be matched.}
% As we refine the target from Black/other
% (\(K=2\), \(c=0.128\)) to ethnicity\(\times\)age
% (\(K=6\), combining three ethnicity groups with two age groups,
% \(c=0.0014\)), the margin \(c\) becomes about \(90\) times smaller.
% At \(T=10^5\)\footnote{Here \(T\) is the number of simulated recruitment
% rounds. The NHIS data are used to estimate the distribution of each
% recruitment source, and recruits are then repeatedly simulated from these
% estimated distributions.}, the product \(c\|S_T\|\) remains of comparable
% size across the different targets: it lies in \([0.07,0.25]\) for the
% oracle and in \([0.12,0.42]\) for
% Algorithm~\ref{alg:ucb-sampling} with radius \(0.3\,r_a\)
% (Appendix Figure~\ref{fig:app-ladder}). Thus, the observed imbalance grows roughly as 
% \(1/c\), consistent with the dependence on \(c\) in
% Theorem~\ref{thm:multidimensional-guarantee}.

% This also makes finer targets harder to match. After \(785\) simulated
% recruits, Algorithm~\ref{alg:ucb-sampling} is \(5.8\) times closer to the
% target than uniform recruitment for ethnicity, but only \(2.7\) times
% closer for ethnicity\(\times\)age. As \(c\) becomes smaller, the available
% recruitment sources give the algorithm less room to correct the sample
% composition.

% \paragraph{Cost-aware sampling.}
% We next ask whether cost-aware sampling buys better representation for a
% fixed budget. In these experiments offline sources cost \(\$108\) per
% recruit, \(1.5\) times the online cost and within the range reported
% by~\citet{brogger2020online}; Appendix~\ref{app:exp-cost} repeats the
% comparison with the median costs. For \(c_0=0.05\), the cheapest
% three-source support with margin at least \(c_0\) is
% \(I^\star=\{\)online, southern clinics, western clinics\(\}\), costing
% \(\$90.70\) per recruit, against \(\$108\) for both the cost-unaware closed
% loop and the myopic rule of Section~\ref{sec:cost}. A feasible mixture
% costing \(\$81.68\) exists, but its margin is only \(0.004\).

% We fix a total budget and compare the error \(\|\widehat p_T-p_G\|_2\),
% averaged over runs, of Algorithm~\ref{alg:cost-support} (\(c_0=0.05\),
% \(n=400\) samples per source, kept in the sample) followed by
% Algorithm~\ref{alg:ucb-sampling}, with cheap \& discard as in Example~1: it
% recruits only online, plans the sample size \(N\) that the budget buys on
% average, keeps each recruit while the quota \(\mathrm{round}(Np_G)\) of their
% group is open, and pays for every contact (Figure~\ref{fig:main}(c)). Below
% about \(\$1\)M, Algorithm~\ref{alg:cost-support} is still exploring
% (\(6{,}400\) recruits costing \(\$662\)k), and cheap \& discard is more
% accurate. From \(\$1.1\)M on, Algorithm~\ref{alg:cost-support}\(\to\)\ref{alg:ucb-sampling}
% is more accurate, by a factor of \(4\) to \(6\) above \(\$1.7\)M; at
% \(\$8\)M its error is \(1.3\times10^{-4}\) against \(6.7\times10^{-4}\).
% The error of cheap \& discard comes from the runs, about half, in which the
% budget runs out before the Hispanic quota is filled.
% Algorithm~\ref{alg:ucb-sampling} on all sources is more accurate still: it
% pays \(18\%\) more per recruit, but needs no exploration phase and cannot
% select a wrong support. The length of exploration matters: with \(n=100\)
% or \(200\), \(31\%\) and \(10\%\) of runs select a support whose hull does
% not contain \(p_G\), and their error stalls (Appendix~\ref{app:exp-cost}).
% \paragraph{Limitations.}
% At small budgets, the exploration phase of Algorithm~\ref{alg:cost-support}
% makes it less accurate than cheap \& discard. With the median costs
% of~\citet{brogger2020online} (offline \(\$199\)),
% Algorithm~\ref{alg:cost-support}\(\to\)\ref{alg:ucb-sampling} overtakes
% cheap \& discard only above about \(\$2.5\)M
% (Appendix~\ref{app:exp-cost}).

% Matching the target demographics also does not guarantee representativeness
% in variables that are not included in the target. For example, after
% matching ethnicity, the recruited sample still overrepresents adults in
% fair or poor health by \(4.1\) percentage points
% (Appendix~\ref{app:exp-downstream}). In practice, variables that strongly
% affect the study endpoint should therefore also be included in the target.
% The tradeoff is that adding more target dimensions generally reduces the
% margin \(c\), making the representation problem harder.

% Overall, our results show that feedback can correct representation error during
% data collection even when source distributions are unknown. The approach extends
% from one group to multiple groups with the optimal \(T^{-1}\) rate, and can be
% combined with support selection when recruitment sources have different costs.

\subsection*{AI use statement}

We used generative AI tools, including ChatGPT and Claude, to assist with
theoretical development, checking and drafting mathematical arguments,
experimental design and implementation, interpretation of results, and
editing the manuscript. All AI-assisted proofs, code, experimental results,
and factual claims were reviewed and verified by the authors. We take
responsibility for the final content of this work, including text, claims,
code, and other artifacts produced with the aid of generative AI.

\bibliographystyle{abbrvnat}
\bibliography{refs}


%%%%%%%%%%%%%%%%%%%%%%%%%%%%%%%%%%%%%%%%%%%%%%%%%%%%%%%%%%%%
\section*{Checklist}

\begin{enumerate}

  \item For all models and algorithms presented, check if you include:
  \begin{enumerate}
    \item A clear description of the mathematical setting, assumptions, algorithm, and/or model.
    [\textbf{Yes}. The mathematical setting and algorithms are described in
    Sections~3--5.]

    \item An analysis of the properties and complexity (time, space, sample size) of any algorithm.
    [\textbf{Yes}. We provide finite-sample and asymptotic guarantees for the
    proposed algorithms.]

    \item (Optional) Anonymized source code, with specification of all dependencies, including external libraries.
    [\textbf{Yes}. Anonymized code and instructions are included in the
    supplementary material.]
  \end{enumerate}

  \item For any theoretical claim, check if you include:
  \begin{enumerate}
    \item Statements of the full set of assumptions of all theoretical results.
    [\textbf{Yes}.]

    \item Complete proofs of all theoretical results.
    [\textbf{Yes}. Complete proofs are provided in the supplementary material.]

    \item Clear explanations of any assumptions.
    [\textbf{Yes}. The main assumptions, including interiority and the source
    assumptions, are discussed in the main text.]
  \end{enumerate}

  \item For all figures and tables that present empirical results, check if you include:
  \begin{enumerate}
    \item The code, data, and instructions needed to reproduce the main experimental results (either in the supplemental material or as a URL).
    [\textbf{Yes}. The experiments use the publicly available 2023 NHIS data,
    and anonymized code and reproduction instructions are included in the
    supplementary material.]

    \item All the training details (e.g., data splits, hyperparameters, how they were chosen).
    [\textbf{Yes}.]

    \item A clear definition of the specific measure or statistics and error bars (e.g., with respect to the random seed after running experiments multiple times).
    [\textbf{Yes}. Experimental protocols, metrics, numbers of replications,
    and variability across runs are described in the supplementary material.]

    \item A description of the computing infrastructure used. (e.g., type of GPUs, internal cluster, or cloud provider).
    [\textbf{Yes}. Experiments were run on an Ubuntu 24.04 VM with 4 vCPUs, 16 GB of RAM, and 30 GB of disk.]
  \end{enumerate}

  \item If you are using existing assets (e.g., code, data, models) or curating/releasing new assets, check if you include:
  \begin{enumerate}
    \item Citations of the creator If your work uses existing assets.
    [\textbf{Yes}. We use and cite the 2023 National Health Interview Survey
    released by the National Center for Health Statistics.]

    \item The license information of the assets, if applicable.
    [\textbf{Yes}. We use the publicly released NHIS files under the NCHS
    public-use data terms.]

    \item New assets either in the supplemental material or as a URL, if applicable.
    [\textbf{Yes}. The anonymized experimental code is included in the
    supplementary material.]

    \item Information about consent from data providers/curators.
    [\textbf{Not Applicable}. We use previously collected public-use survey
    data and did not recruit participants or collect new human-subject data.]

    \item Discussion of sensible content if applicable, e.g., personally identifiable information or offensive content.
    [\textbf{Yes}. We use only the public-use NHIS files, from which
    potentially identifying information has been removed.]
  \end{enumerate}

  \item If you used crowdsourcing or conducted research with human subjects, check if you include:
  \begin{enumerate}
    \item The full text of instructions given to participants and screenshots.
    [\textbf{Not Applicable}.]

    \item Descriptions of potential participant risks, with links to Institutional Review Board (IRB) approvals if applicable.
    [\textbf{Not Applicable}.]

    \item The estimated hourly wage paid to participants and the total amount spent on participant compensation.
    [\textbf{Not Applicable}.]
  \end{enumerate}

\end{enumerate}
\clearpage
\appendix
\thispagestyle{empty}
\onecolumn

\section{Proofs}
\begin{proof}[of Theorem \ref{thm:one-dimensional-guarantee} ]
Let \(\mathcal F_t\) denote the history of the sampling process up to time
\(t\). Our goal is to show that \(S_T\) remains bounded with high probability,
uniformly in the horizon \(T\). A convenient way to obtain such a tail bound
is to control an exponential moment of \(|S_T|\). Indeed, if for some
\(\lambda>0\) we can show that
\[
\sup_{T\geq0}
\mathbb E\left[e^{\lambda |S_T|}\right]
<\infty,
\]
then Markov's inequality immediately gives an exponentially decaying bound on
\(\mathbb P(|S_T|\geq k)\). We therefore consider the exponential potential
\[
V_t:=e^{\lambda |S_t|}
\]
and show that it contracts in expectation whenever \(S_t\) is away from zero.

Recall that
\[
S_{t+1}
=
S_t+Y_{t+1}-p_G.
\]
Under phone sampling,
\[
\mathbb E[Y_{t+1}-p_G\mid A_{t+1}=P]
=
p_P-p_G<0,
\]
whereas under online sampling,
\[
\mathbb E[Y_{t+1}-p_G\mid A_{t+1}=O]
=
p_O-p_G>0.
\]

Suppose first that \(S_t>1\). Since
\[
Y_{t+1}-p_G\in\{-p_G,1-p_G\}
\]
and both possible increments have absolute value less than \(1\),
\(S_{t+1}\) cannot cross zero in one step. Hence
\[
|S_{t+1}|
=
S_t+Y_{t+1}-p_G.
\]
The algorithm chooses phone sampling when \(S_t>0\), and therefore
\[
\begin{aligned}
\mathbb E[V_{t+1}\mid\mathcal F_t]
&=
e^{\lambda S_t}
\mathbb E\left[
e^{\lambda(Y_{t+1}-p_G)}
\mid\mathcal F_t
\right] \\
&=
V_t\,g_P(\lambda),
\end{aligned}
\]
where
\[
g_P(\lambda)
:=
p_Pe^{\lambda(1-p_G)}
+
(1-p_P)e^{-\lambda p_G}.
\]
At \(\lambda=0\),
\[
g_P(0)=1,
\]
while
\[
g_P'(0)
=
p_P-p_G<0.
\]
Hence, by continuity, for all sufficiently small \(\lambda>0\),
\[
g_P(\lambda)<1.
\]

Now suppose that \(S_t<-1\). Again \(S_{t+1}\) cannot cross zero in one
step, and therefore
\[
|S_{t+1}|
=
-\bigl(S_t+Y_{t+1}-p_G\bigr).
\]
In this case the algorithm chooses online sampling. Thus
\[
\begin{aligned}
\mathbb E[V_{t+1}\mid\mathcal F_t]
&=
e^{-\lambda S_t}
\mathbb E\left[
e^{-\lambda(Y_{t+1}-p_G)}
\mid\mathcal F_t
\right] \\
&=
V_t\,g_O(\lambda),
\end{aligned}
\]
where
\[
g_O(\lambda)
:=
p_Oe^{-\lambda(1-p_G)}
+
(1-p_O)e^{\lambda p_G}.
\]
Again,
\[
g_O(0)=1,
\]
and
\[
g_O'(0)
=
p_G-p_O<0.
\]
Thus, for all sufficiently small \(\lambda>0\),
\[
g_O(\lambda)<1.
\]

We may therefore choose a single \(\lambda>0\) such that both inequalities
hold, and define
\[
\rho
:=
\max\{g_P(\lambda),g_O(\lambda)\}<1.
\]
We have shown that whenever \(|S_t|>1\),
\[
\mathbb E[V_{t+1}\mid\mathcal F_t]
\leq
\rho V_t.
\]

It remains to consider the region \(|S_t|\leq1\). Since
\[
|Y_{t+1}-p_G|\leq1,
\]
we have
\[
|S_{t+1}|
\leq
|S_t|+|Y_{t+1}-p_G|
\leq2.
\]
Consequently,
\[
V_{t+1}
=
e^{\lambda|S_{t+1}|}
\leq
e^{2\lambda}.
\]

Combining the two regions gives the one-step conditional inequality
\[
\mathbb E[V_{t+1}\mid\mathcal F_t]
\leq
\rho V_t+e^{2\lambda}.
\]
This inequality should be understood conditionally on the history up to time
\(t\). Once \(\mathcal F_t\) is fixed, the current deficit \(S_t\), the
potential \(V_t\), and the source selected by the algorithm are all known; the
only remaining randomness comes from the next respondent. Thus, for each
realized history, the inequality bounds the expected value of the potential
after one further sample. Since \(V_t\) depends on the realized history, the
right-hand side is itself random before conditioning on \(\mathcal F_t\).

Taking expectations of both sides and using the tower property,
\[
\begin{aligned}
\mathbb E[V_{t+1}]
&=
\mathbb E\left[
\mathbb E[V_{t+1}\mid\mathcal F_t]
\right] \\
&\leq
\rho\,\mathbb E[V_t]
+
e^{2\lambda}.
\end{aligned}
\]
Since \(S_0=0\),
\[
V_0=1.
\]
Iterating the preceding recursion gives
\[
\begin{aligned}
\mathbb E[V_T]
&\leq
\rho^T V_0
+
e^{2\lambda}
\sum_{j=0}^{T-1}\rho^j \\
&\leq
1+\frac{e^{2\lambda}}{1-\rho}.
\end{aligned}
\]
Define
\[
C
:=
1+\frac{e^{2\lambda}}{1-\rho}.
\]
Then
\[
\sup_{T\geq0}
\mathbb E\left[e^{\lambda|S_T|}\right]
\leq C.
\]

We can now convert this exponential-moment bound into a tail bound. By
Markov's inequality, for every \(k\geq0\),
\[
\begin{aligned}
\mathbb P(|S_T|\geq k)
&=
\mathbb P\left(
e^{\lambda|S_T|}
\geq
e^{\lambda k}
\right) \\
&\leq
e^{-\lambda k}
\mathbb E\left[e^{\lambda|S_T|}\right] \\
&\leq
Ce^{-\lambda k}.
\end{aligned}
\]
This proves the first assertion.

Finally, given \(\delta\in(0,1)\), take
\[
k
=
\frac{1}{\lambda}
\log\left(\frac{C}{\delta}\right).
\]
Then
\[
\mathbb P\left(
|S_T|
\geq
\frac{1}{\lambda}
\log\left(\frac{C}{\delta}\right)
\right)
\leq
\delta.
\]
Since
\[
S_T
=
T(\widehat p_T-p_G),
\]
it follows that, with probability at least \(1-\delta\),
\[
\left|
\widehat p_T-p_G
\right|
\leq
\frac{1}{\lambda T}
\log\left(\frac{C}{\delta}\right).
\]
This completes the proof.
\end{proof}
\subsection{Proof of Theorem~\ref{thm:multidimensional-guarantee}.}

The proof uses the concept of Upper Confidence Bounds (UCB), which is commonly
used in multi-armed bandit problems. At every round of sampling, we can measure
our current error \(S_t\) relative to the target distribution. What we do not
know is the distribution associated with each source. If we knew these
distributions, it would be much easier to choose a source that would reduce our
sampling error. However, with each round of sampling, we obtain more information
about the sources that have been sampled so far. This helps us choose a better
candidate for the next round and, eventually, keep the final sampling error
bounded.

The proof will proceed as follows. We assign a confidence radius \(r_a(n)\) to
each source. This tells us how close the empirical mean of source \(a\), after \(n\)
samples, is to its actual mean. If the confidence radius is smaller than
\(c/4\),\footnote{Recall that \(c\) is the radius of the ball centered at
\(p_G\) that lies inside the convex hull of the source means \(\mu_a\).} then we
can choose a source that reduces our sampling error. Therefore, our algorithm will
encourage sampling sources that are still poorly known. As more samples are
collected from each source, its confidence radius becomes smaller, allowing the
algorithm to make increasingly reliable choices of which source to sample next.

There are several subtleties in the proof, which we will mention and discuss as
they arise. One important issue is that the confidence bounds themselves are
probabilistic statements. For instance, it is possible that the empirical mean
of one of the sources remains far from its actual mean for a certain number of
samples. Therefore, we first need to show that, with high probability, all of
the confidence bounds used by the algorithm are simultaneously valid. We begin
by taking care of this issue.


For each \(t\geq 0\), define the event

$$
\mathcal E_t
:=
\bigcap_{a=1}^m
\bigcap_{1\le n\le N_a(t)}
\left\{
\|\widehat\mu_{a,n}-\mu_a\|
\le
r_a(n)
\right\}.
$$

Thus, \(\mathcal E_t\) is the event that every confidence bound corresponding
to observations available by time \(t\) is valid. In particular,  $
\mathcal E_t\in\mathcal F_t,
$
and the events are decreasing: 
$
\mathcal E_{t+1}\subseteq\mathcal E_t.
$
Define the all-time confidence event
\[
\mathcal E
:=
\bigcap_{t\ge0}\mathcal E_t.
\]
Thus, \(\mathcal E\) is the event that every confidence bound actually used
by the algorithm is valid throughout the sampling process.

We choose the confidence radii so that, for every source \(a\) and every
\(n\ge1\),
\begin{equation}
\label{eqdmn}    \mathbb P\left(
    \|\widehat\mu_{a,n}-\mu_a\|>r_a(n)
\right)
\le
\delta_{a,n},
\end{equation}
where
\begin{equation}
\label{eqd}
\sum_{a=1}^m\sum_{n\ge1}\delta_{a,n}\le\delta.
\end{equation}

Therefore, by a union bound,
\[
\mathbb P(\mathcal E^c)
\le
\sum_{a=1}^m\sum_{n\ge1}\delta_{a,n}
\le
\delta,
\]
and hence $
\mathbb P(\mathcal E)\ge1-\delta.
$
For \(n\ge1\), define
\[
    \delta_{a,n}
    :=
    \frac{6\delta}{\pi^2 m n^2},
\]
and choose the confidence radius
\[
    r_a(n)
    :=
    \sqrt{
        \frac{K}{2n}
        \log\left(
            \frac{\pi^2 m K n^2}{3\delta}
        \right)
    }.
\]
For a source that has not yet been sampled, we set $  r_a(0):=+\infty.$ These choices ensure that \eqref{eqdmn} and \eqref{eqd},
and that
\[
    r_a(n)\longrightarrow0
    \qquad\text{as }n\to\infty.
\]
These facts follow directly from Hoeffding's inequality and a union bound
over all coordinates; we leave the straightforward verification to the
reader.\\

We therefore analyze the algorithm through the events \(\mathcal E_t\).
Whenever \(\mathcal E_t\) holds, all confidence bounds used by the algorithm
at time \(t\) are valid. The only purpose of keeping the time-indexed events
\(\mathcal E_t\), rather than conditioning directly on the all-time event
\(\mathcal E\), is that \(\mathcal E_t\) is determined by the history
\(\mathcal F_t\), whereas \(\mathcal E\) also contains information about
future observations.

    

% the quantity that enters the expectation argument is
% \[
%     W_t\mathbf 1_{\mathcal E_t}.
% \]
% Since
% \[
%     \mathcal E=\bigcap_{t\ge0}\mathcal E_t
%     \qquad\text{and}\qquad
%     \mathbb P(\mathcal E)\ge1-\delta,
% \]
% with probability at least \(1-\delta\), all of these indicators are equal to
% one throughout the sampling process. \\

We call round \(t+1\) \emph{uncertain} if
\[
    r_{A_{t+1}}\!\left(N_{A_{t+1}}(t)\right)
    >
    \frac{c}{4}.
\]
Otherwise, we call the round \emph{good}. Thus, on a good round the selected
source is known accurately enough for us to guarantee that its true mean points
toward the target distribution. On an uncertain round we may not yet have such
a guarantee, but each time such a source is selected we obtain another sample
from it and hence reduce its uncertainty.
Let \(b_t\) denote the number of uncertain rounds among the first \(t\).
% rounds, and define
% \[
%     W_t
%     :=
%     \exp\left(
%         \lambda\|S_t\|-\alpha b_t
%     \right),
% \]
% where \(\lambda>0\) and \(\alpha>0\) will be specified later.

Main part of the proof is a recursive inequlity for 
\begin{equation}
\label{Wt}
    W_t
    :=
    \exp\left(
        \lambda\|S_t\|-\alpha b_t
    \right),
\end{equation}

Suppose that we have established constants \(0<\rho<1\) and \(B<\infty\)
such that, whenever \(\mathcal E_t\) holds,
\begin{equation}
    \mathbb E\!\left[
        W_{t+1}\mid\mathcal F_t
    \right]
    \le
    \rho W_t+B.
    \tag{R5.20}
    \label{eq:R5.20}
\end{equation}
Since
\(\mathcal E_{t+1}\subseteq\mathcal E_t\),
\[
\mathbb E\left[
    W_{t+1}\mathbf 1_{\mathcal E_{t+1}}
\right]
\le
\mathbb E\left[
    W_{t+1}\mathbf 1_{\mathcal E_t}
\right].
\]
We now apply the tower property:
\[
\mathbb E\left[
    W_{t+1}\mathbf 1_{\mathcal E_t}
\right]
=
\mathbb E\left[
    \mathbb E\left[
        W_{t+1}\mathbf 1_{\mathcal E_t}
        \mid\mathcal F_t
    \right]
\right].
\]
Since \(\mathcal E_t\in\mathcal F_t\), once the history
\(\mathcal F_t\) is fixed, the indicator
\(\mathbf 1_{\mathcal E_t}\) is fixed as well. Hence it can be pulled
outside the conditional expectation:
\[
\mathbb E\left[
    W_{t+1}\mathbf 1_{\mathcal E_t}
    \mid\mathcal F_t
\right]
=
\mathbf 1_{\mathcal E_t}
\mathbb E\left[
    W_{t+1}\mid\mathcal F_t
\right].
\]
Therefore,
\[
\mathbb E\left[
    W_{t+1}\mathbf 1_{\mathcal E_t}
\right]
=
\mathbb E\left[
    \mathbf 1_{\mathcal E_t}
    \mathbb E\left[
        W_{t+1}\mid\mathcal F_t
    \right]
\right].
\]
Using \eqref{eq:R5.20} on the event \(\mathcal E_t\), we obtain
\[
\begin{aligned}
\mathbb E\left[
    W_{t+1}\mathbf 1_{\mathcal E_{t+1}}
\right]
&\le
\rho\,
\mathbb E\left[
    W_t\mathbf 1_{\mathcal E_t}
\right]
+
B.
\end{aligned}
\]


\noindent Since \(S_0=0\), \(b_0=0\), and \(\mathcal E_0\) is the whole sample space, 
we have $W_0\mathbf 1_{\mathcal E_0}=1.$ 
Iterating the preceding inequality gives
\[
\begin{aligned}
    \mathbb E\left[
        W_T\mathbf 1_{\mathcal E_T}
    \right]
    &\le
    \rho^T
    +
    B\sum_{j=0}^{T-1}\rho^j
    \\
    &\le
    1+\frac{B}{1-\rho}.
\end{aligned}
\]
Hence, defining
\begin{equation}
    C
    :=
    1+\frac{B}{1-\rho},
    \tag{R5.21}
    \label{eq:R5.21}
\end{equation}
we obtain
\begin{equation}
    \boxed{
        \mathbb E\left[
            W_T\mathbf 1_{\mathcal E_T}
        \right]
        \le C
    }
    \qquad
    \text{for every }T.
    \tag{R5.22}
    \label{eq:R5.22}
\end{equation}

We will prove below, in Lemma~\ref{lem:finite-uncertain-rounds}, that there
exists a finite constant \(M(c)\), independent of \(T\), such that the total
number of uncertain rounds satisfies
\begin{equation}
    b_T\le M(c)
    \qquad\text{for every }T.
    \tag{R5.23}
    \label{eq:R5.23}
\end{equation}

Applying Lemma~\ref{lem:potential-to-representation} to
\[
    W_T
    =
    \exp\left(
        \lambda\|S_T\|-\alpha b_T
    \right)
\]

and using \eqref{eq:R5.22} and \eqref{eq:R5.23}, we conclude that, for every
\(\eta\in(0,1)\),\footnote{\(\eta\) is the additional failure probability
introduced when applying Markov's inequality to the potential bound.}
with probability at least \(1-\delta-\eta\),

\begin{equation}
    \|S_T\|
    \le
    \frac{\alpha}{\lambda}M(c)
    +
    \frac{1}{\lambda}
    \log\frac{C}{\eta}.
    \tag{R5.24}
    \label{eq:R5.24}
\end{equation}
Here, the additional failure probability \(\delta\) comes from the complement
of the all-time confidence event \(\mathcal E\).

Finally, since
\[
    S_T
    =
    T(\widehat p_T-p_G),
\]
we obtain
\begin{equation}
    \boxed{
        \|\widehat p_T-p_G\|
        \le
        \frac{1}{T}
        \left[
            \frac{\alpha}{\lambda}M(c)
            +
            \frac{1}{\lambda}
            \log\frac{C}{\eta}
        \right].
    }
    \tag{R5.25}
    \label{eq:R5.25}
\end{equation}

Thus, once the recursive bound \eqref{eq:R5.20} is established with constants
\(\lambda,\alpha,\rho,B\) independent of \(T\), the desired
\(O(T^{-1})\) representation guarantee follows.

\begin{lemma}[From a bounded potential to representation error]
\label{lem:potential-to-representation}
Fix \(T\ge1\), and suppose
\[
    W_T
    :=
    \exp\left(
        \lambda\|S_T\|-\alpha b_T
    \right),
\]
where \(\lambda>0\), \(\alpha>0\), and \(b_T\) is the number of uncertain
rounds up to time \(T\). Under the confidence-event convention introduced
above, assume that
\[
\mathbb E\!\left[
W_T\mathbf 1_{\mathcal E_T}
\right]
\le C
\qquad\text{and}\qquad
b_T\le M(c)
\]
almost surely.

Then, for every \(\eta\in(0,1)\), with probability at least
\(1-\delta-\eta\),
\begin{equation}
    \|\widehat p_T-p_G\|
    \le
    \frac{1}{T}
    \left[
        \frac{\alpha}{\lambda}M(c)
        +
        \frac{1}{\lambda}
        \log\frac{C}{\eta}
    \right].
    \label{eq:potential-to-representation-bound}
\end{equation}
\end{lemma}

\paragraph{Proof of the recursive inequality.}


It remains to establish \eqref{eq:R5.20}. Fix a time \(t\) and suppose that
\(\mathcal E_t\) holds. For \(S_t\neq0\), write
\[
    u_t:=\frac{S_t}{\|S_t\|}.
\]
Recall that at round \(t+1\), the algorithm selects
\[
    A_{t+1}
    \in
    \arg\min_{a}
    \left\{
        \left\langle
            u_t,
            \widehat\mu_{a,N_a(t)}-p_G
        \right\rangle
        -
        r_a\!\left(N_a(t)\right)
    \right\}.
\]
We distinguish three
cases.

\medskip

\noindent
\textbf{Case 1: the round is uncertain.} On an uncertain round,
\[
    b_{t+1}=b_t+1.
\]
Moreover,
\[
    S_{t+1}
    =
    S_t+(Y_{t+1}-p_G),
\]
and since
\[
    \|Y_{t+1}-p_G\|\le R,
\]
the triangle inequality gives
\[
    \|S_{t+1}\|
    \le
    \|S_t\|+R.
\]
Consequently,
\[
\begin{aligned}
    W_{t+1}
    &=
    \exp\left(
        \lambda\|S_{t+1}\|-\alpha b_{t+1}
    \right)
    \\
    &\le
    \exp\left(
        \lambda(\|S_t\|+R)-\alpha(b_t+1)
    \right)
    \\
    &=
    W_t e^{\lambda R-\alpha}.
\end{aligned}
\]
We choose
\begin{equation}
    \alpha
    :=
    \lambda R+\kappa,
    \tag{R5.26}
    \label{eq:R5.26}
\end{equation}
where \(\kappa>0\) will be specified below. Then
\begin{equation}
    W_{t+1}
    \le
    e^{-\kappa}W_t.
    \tag{R5.27}
    \label{eq:R5.27}
\end{equation}

Thus, although an uncertain round may increase \(\|S_t\|\), the additional
penalty in \(b_t\) compensates for this possible increase.

\medskip

\noindent
\textbf{Case 2: the round is good and \(\|S_t\|\) is sufficiently large.} 
On a good round, the mean of the selected source points toward the target, but the
actual observation is still random, so the realized error need not decrease on
every round. We therefore first work with the squared norm, because expanding
\(\|S_{t+1}\|^2\) exposes the directional drift through an inner product.
We then use Jensen's inequality to convert the resulting bound on the expected
squared norm into a bound on the expected norm itself. Finally, since our
potential is exponential in \(\|S_t\|\), we use Hoeffding's lemma to turn this
negative expected change in the norm, together with the bounded size of one-step
fluctuations, into a contraction of the exponential potential.\\

In the categorical case, both \(Y_t\) and \(p_G\) have coordinates summing to one. Hence the coordinates of \(S_t\) sum to zero, so
\(u_t=S_t/\|S_t\|\) lies in the direction space of the simplex. Since \(p_G\) has relative-interiority margin \(c\), there exists a source \(a\) such that
\[
\langle u_t,\mu_a-p_G\rangle\le -c.
\]

More precisely, by the interiority assumption, there exists a source whose mean has directional drift at most \(-c\); on a good round, the confidence bounds and the UCB selection rule imply that the selected source still has directional drift at most $-c/2$
\[
    \left\langle
        \frac{S_t}{\|S_t\|},
        \mu_{A_{t+1}}-p_G
    \right\rangle
    \le
    -\frac c2.
\]
Set $\gamma:=\frac c2.$ Since
\[
    \mathbb E[
        Y_{t+1}-p_G
        \mid\mathcal F_t
    ]
    =
    \mu_{A_{t+1}}-p_G,
\]
it follows that
\begin{equation}
    \left\langle
        S_t,
        \mathbb E[
            Y_{t+1}-p_G
            \mid\mathcal F_t
        ]
    \right\rangle
    \le
    -\gamma\|S_t\|.
    \tag{R5.28}
    \label{eq:R5.28}
\end{equation}

Using
\[
    S_{t+1}
    =
    S_t+(Y_{t+1}-p_G),
\]
we obtain
\[
\begin{aligned}
    \mathbb E[
        \|S_{t+1}\|^2
        \mid\mathcal F_t
    ]
    &=
    \|S_t\|^2
    +
    2
    \left\langle
        S_t,
        \mathbb E[
            Y_{t+1}-p_G
            \mid\mathcal F_t
        ]
    \right\rangle
    \\
    &\qquad
    +
    \mathbb E[
        \|Y_{t+1}-p_G\|^2
        \mid\mathcal F_t
    ].
\end{aligned}
\]
Since
\[
    \|Y_{t+1}-p_G\|\le R,
\]
\eqref{eq:R5.28} gives
\[
    \mathbb E[
        \|S_{t+1}\|^2
        \mid\mathcal F_t
    ]
    \le
    \|S_t\|^2
    -
    2\gamma\|S_t\|
    +
    R^2.
\]

Let
\begin{equation}
    L
    :=
    \max\left\{
        \frac{R^2}{\gamma},
        \gamma
    \right\}.
    \tag{R5.29}
    \label{eq:R5.29}
\end{equation}
Whenever \(\|S_t\|\ge L\),
\[
    R^2
    \le
    \gamma\|S_t\|,
\]
and therefore
\[
    \mathbb E[
        \|S_{t+1}\|^2
        \mid\mathcal F_t
    ]
    \le
    \|S_t\|^2-\gamma\|S_t\|.
\]
By Jensen's inequality,
\[
\begin{aligned}
    \mathbb E[
        \|S_{t+1}\|
        \mid\mathcal F_t
    ]
    &\le
    \sqrt{
        \mathbb E[
            \|S_{t+1}\|^2
            \mid\mathcal F_t
        ]
    }
    \\
    &\le
    \sqrt{
        \|S_t\|^2-\gamma\|S_t\|
    }.
\end{aligned}
\]
Since $\|S_t\|\ge L\ge\gamma,$
we have
\[
    \sqrt{
        \|S_t\|^2-\gamma\|S_t\|
    }
    \le
    \|S_t\|-\frac{\gamma}{2}.
\]
Thus
\begin{equation}
    \mathbb E\left[
        \|S_{t+1}\|-\|S_t\|
        \mid\mathcal F_t
    \right]
    \le
    -\frac{\gamma}{2}.
    \tag{R5.30}
    \label{eq:R5.30}
\end{equation}

Now set
\[
    \Delta_t
    :=
    \|S_{t+1}\|-\|S_t\|.
\]
By the reverse triangle inequality,
\[
    |\Delta_t|
    \le
    \|Y_{t+1}-p_G\|
    \le
    R.
\]
Therefore, conditionally on \(\mathcal F_t\), the random variable
\(\Delta_t\) lies in an interval of length at most \(2R\). Hoeffding's
lemma and \eqref{eq:R5.30} give
\[
\begin{aligned}
    \mathbb E[
        e^{\lambda\Delta_t}
        \mid\mathcal F_t
    ]
    &\le
    \exp\left(
        \lambda
        \mathbb E[
            \Delta_t
            \mid\mathcal F_t
        ]
        +
        \frac{\lambda^2R^2}{2}
    \right)
    \\
    &\le
    \exp\left(
        -\frac{\lambda\gamma}{2}
        +
        \frac{\lambda^2R^2}{2}
    \right).
\end{aligned}
\]

Choose
\begin{equation}
    \lambda
    :=
    \frac{\gamma}{2R^2}.
    \tag{R5.31}
    \label{eq:R5.31}
\end{equation}
Then
\[
    -\frac{\lambda\gamma}{2}
    +
    \frac{\lambda^2R^2}{2}
    =
    -\frac{\gamma^2}{8R^2}.
\]
We therefore set
\begin{equation}
    \kappa
    :=
    \frac{\gamma^2}{8R^2},
    \qquad
    \rho
    :=
    e^{-\kappa}<1.
    \tag{R5.32}
    \label{eq:R5.32}
\end{equation}
It follows that
\[
    \mathbb E[
        e^{\lambda\|S_{t+1}\|}
        \mid\mathcal F_t
    ]
    \le
    \rho
    e^{\lambda\|S_t\|}.
\]
Since a good round does not increase \(b_t\),
\begin{equation}
    \mathbb E[
        W_{t+1}
        \mid\mathcal F_t
    ]
    \le
    \rho W_t
    \qquad
    \text{if the round is good and }\|S_t\|\ge L.
    \tag{R5.33}
    \label{eq:R5.33}
\end{equation}
Notice also that, by \eqref{eq:R5.26},
\[
    e^{-\kappa}
    =
    \rho,
\]
so \eqref{eq:R5.27} gives
\[
    W_{t+1}
    \le
    \rho W_t
\]
on every uncertain round as well.

\medskip

\noindent
\textbf{Case 3: the round is good and \(\|S_t\|<L\).} In this case,
\[
    \|S_{t+1}\|
    \le
    \|S_t\|+R
    <
    L+R.
\]
Since \(b_{t+1}=b_t\) on a good round and \(b_t\ge0\),
\[
    W_{t+1}
    \le
    e^{\lambda(L+R)}.
\]
Define
\begin{equation}
    B
    :=
    e^{\lambda(L+R)}.
    \tag{R5.34}
    \label{eq:R5.34}
\end{equation}
Then
\begin{equation}
    W_{t+1}
    \le
    B
    \qquad
    \text{if the round is good and }\|S_t\|<L.
    \tag{R5.35}
    \label{eq:R5.35}
\end{equation}

Combining the three cases, whenever \(\mathcal E_t\) holds,
\begin{equation}
    \mathbb E[
        W_{t+1}\mid\mathcal F_t
    ]
    \le
    \rho W_t+B.
    \tag{R5.36}
    \label{eq:R5.36}
\end{equation}
Indeed, Cases 1 and 2 give contraction by the factor \(\rho\), while
Case 3 gives the uniform bound \(B\). Under the confidence-event convention
established above, this is precisely the recursive inequality
\eqref{eq:R5.20}.

Finally, since
\[
    \gamma=\frac c2,
    \qquad
    \lambda
    =
    \frac{\gamma}{2R^2}
    =
    \frac{c}{4R^2},
\]
we have
\[
    \kappa
    =
    \frac{\gamma^2}{8R^2}
    =
    \frac{c^2}{32R^2},
    \qquad
    \rho
    =
    \exp\left(
        -\frac{c^2}{32R^2}
    \right),
\]
and
\[
    \frac{\alpha}{\lambda}
    =
    R+\frac{\kappa}{\lambda}
    =
    R+\frac c8,
    \qquad
    \frac1\lambda
    =
    \frac{4R^2}{c}.
\]
Moreover,
\[
    B
    =
    e^{\lambda(L+R)}.
\]
All of these constants are independent of \(T\). Thus the recursive bound
assumed earlier has been established, and the bounded-potential argument
completes the proof.

\begin{lemma}[Finite number of uncertain rounds]
\label{lem:finite-uncertain-rounds}
For the confidence radius \(r_a(n)\) defined above, there exists a finite
integer \(N(c)\) such that
\[
    r_a(n)\le \frac{c}{4}
    \qquad\text{for every source }a
    \text{ and every }n\ge N(c).
\]
Consequently,
\[
    b_T\le M(c):=mN(c)
    \qquad\text{for every }T.
\]
\end{lemma}

\begin{proof}
Since
\[
    r_a(n)
    =
    \sqrt{
        \frac{K}{2n}
        \log\left(
            \frac{\pi^2mKn^2}{3\delta}
        \right)
    }
    \longrightarrow 0,
\]
there exists a finite \(N(c)\) such that
\[
    r_a(n)\le\frac{c}{4}
    \qquad
    \text{for all }n\ge N(c).
\]
Thus, a source can be selected on an uncertain round only while it has been
sampled fewer than \(N(c)\) times. Each source can therefore generate at most
\(N(c)\) uncertain rounds. Summing over the \(m\) sources gives
\[
    b_T\le mN(c)=M(c)
\]
for every \(T\).
\end{proof}

\begin{proof}
Under the confidence-event convention introduced above,
\[
    \mathbb E\left[
        W_T\mathbf 1_{\mathcal E_T}
    \right]
    \le C.
\]
Therefore, by Markov's inequality,
\[
\begin{aligned}
    \mathbb P\left(
        \mathcal E_T
        \cap
        \left\{
            W_T>\frac{C}{\eta}
        \right\}
    \right)
    &\le
    \frac{\eta}{C}
    \mathbb E\left[
        W_T\mathbf 1_{\mathcal E_T}
    \right]
    \\
    &\le
    \eta.
\end{aligned}
\]
Since
\[
    \mathcal E\subseteq\mathcal E_T
    \qquad\text{and}\qquad
    \mathbb P(\mathcal E^c)\le\delta,
\]
with probability at least \(1-\delta-\eta\), the event \(\mathcal E\) holds
and
\[
    W_T
    \le
    \frac{C}{\eta}.
\]

Using the definition of \(W_T\), this implies
\[
    \lambda\|S_T\|-\alpha b_T
    \le
    \log\frac{C}{\eta}.
\]
Therefore,
\[
    \|S_T\|
    \le
    \frac{\alpha}{\lambda}b_T
    +
    \frac{1}{\lambda}
    \log\frac{C}{\eta}.
\]
Since \(b_T\le M(c)\),
\[
    \|S_T\|
    \le
    \frac{\alpha}{\lambda}M(c)
    +
    \frac{1}{\lambda}
    \log\frac{C}{\eta}.
\]
Finally,
\[
    S_T
    =
    T(\widehat p_T-p_G),
\]
and dividing by \(T\) gives
\eqref{eq:potential-to-representation-bound}.
\end{proof}
\subsection{Proof of the Cost Theorem}
The proof follows the structure of Algorithm~\ref{alg:cost-support}. We first
show that, if each source is explored sufficiently many times, then the empirical
source means are accurate enough to determine which collections of sources have a
usable interior margin around \(p_G\). In particular,
Algorithm~\ref{alg:cost-support} keeps every support that is truly good enough,
and any support it selects still has enough margin around \(p_G\).

We then show that the same exploration accuracy is enough to estimate the cost
associated with each admissible support. Therefore
Algorithm~\ref{alg:cost-support}, which selects the support with the smallest
empirical cost, returns a collection of sources whose true cost is close to the
cheapest feasible one. Finally, once this support has been selected, we
restrict the multidimensional adaptive sampling algorithm to those sources. The
representation theorem then gives the desired convergence to \(p_G\), while
the support-selection argument gives the near-minimal cost guarantee.

\begin{proof}[Proof of Theorem~\ref{thm:cost-support}]
Let
\[
\varepsilon_n
:=
\sqrt{
\frac{K}{2n}
\log\frac{2mK}{\delta}
}.
\]
By the lower bound on \(n\) assumed in the theorem,
\[
\varepsilon_n
\leq
\frac{c_0}{4}.
\]

Now define
\[
\mathcal E
:=
\left\{
\max_{a\in[m]}
\|\widehat\mu_a-\mu_a\|_2
\leq
\varepsilon_n
\right\}.
\]
By Hoeffding's inequality and a union bound,
\[
\mathbb P(\mathcal E)\geq 1-\delta.
\]
We work on \(\mathcal E\).
If a support \(I\) is considered by Algorithm~\ref{alg:cost-support}, then
\[
\widehat c(I)\geq \frac{3c_0}{4}.
\]
By Lemma~\ref{lem:margin-stability} and
\(\varepsilon_n\leq c_0/4\),
\[
c(I)
\geq
\widehat c(I)-\varepsilon_n
\geq
\frac{c_0}{2}.
\]

Now let
\[
I^\star
\in
\operatorname*{arg\,min}_{\substack{|I|=K\\c(I)\geq c_0}}
B(I).
\]
Every support with true margin at least \(c_0\) is included in the
empirical optimization set in Algorithm~\ref{alg:cost-support}. Indeed, if \(c(I)\geq c_0\), then
Lemma~\ref{lem:margin-stability} gives
\[
\widehat c(I)
\geq
c(I)-\varepsilon_n
\geq
c_0-\frac{c_0}{4}
=
\frac{3c_0}{4}.
\]
Thus
\[
\left\{
I:\ |I|=K,\ c(I)\geq c_0
\right\}
\subseteq
\left\{
I:\ |I|=K,\ \widehat c(I)\geq\frac{3c_0}{4}
\right\}.
\]

The empirical optimization set may be strictly larger. For example, a
support with
\[
c(I)=0.6c_0
\]
could have
\[
\widehat c(I)=0.8c_0,
\]
which is consistent with
\[
|\widehat c(I)-c(I)|
\leq
\frac{c_0}{4},
\]
and would therefore also be included in the search.

Therefore, \(I^\star\) is included in the empirical optimization set.
Since \(\widehat I\) is the support returned by
Algorithm~\ref{alg:cost-support} and, by definition, minimizes
\(\widehat B(I)\) over this possibly larger set,
\[
\widehat B(\widehat I)
\leq
\widehat B(I^\star).
\]

By Lemma~\ref{lem:cost-stability}, since
\(c(\widehat I)\geq c_0/2\),
\[
|\widehat B(\widehat I)-B(\widehat I)|
\leq
\frac{2\ell_{\max}}{c_0}\varepsilon_n.
\]
Likewise, since \(c(I^\star)\geq c_0\),
\[
|\widehat B(I^\star)-B(I^\star)|
\leq
\frac{\ell_{\max}}{c_0}\varepsilon_n.
\]
Therefore
\[
\begin{aligned}
B(\widehat I)
&\leq
\widehat B(\widehat I)
+
\frac{2\ell_{\max}}{c_0}\varepsilon_n\\
&\leq
\widehat B(I^\star)
+
\frac{2\ell_{\max}}{c_0}\varepsilon_n\\
&\leq
B(I^\star)
+
\frac{3\ell_{\max}}{c_0}\varepsilon_n.
\end{aligned}
\]
Recall that
\[
B(I^\star)
=
\min_{\substack{|I|=K\\c(I)\geq c_0}}B(I),
\]
this proves the claimed bound on \(\mathcal E\), and hence with
probability at least \(1-\delta\).
\end{proof}

\begin{lemma}[Stability of the margin]
\label{lem:margin-stability}
Let \(I\subseteq[m]\), and suppose
\[
\max_{a\in I}
\|\widehat\mu_a-\mu_a\|_2
\leq
\varepsilon.
\]
Then
\[
\bigl|
\widehat c(I)-c(I)
\bigr|
\leq
\varepsilon.
\]
\end{lemma}

\begin{proof}[Proof of Lemma~\ref{lem:margin-stability}]
Each estimated source mean \(\widehat\mu_a\) is within distance
\(\varepsilon\) of the corresponding true mean \(\mu_a\). Therefore every
supporting boundary of the empirical convex hull can move by at most
\(\varepsilon\) from the corresponding boundary of the true convex hull:
in any direction, moving each vertex by at most \(\varepsilon\) changes its
projection in that direction by at most \(\varepsilon\).

Since \(c(I)\) is the distance from \(p_G\) to the nearest boundary of the
true convex hull, and \(\widehat c(I)\) is the corresponding distance for the
empirical convex hull, these distances can differ by at most
\(\varepsilon\). Hence
\[
|\widehat c(I)-c(I)|
\leq
\varepsilon.
\]
\end{proof}

% \begin{proof}[Proof of Lemma~\ref{lem:margin-stability}]
% The idea is that perturbing every source mean by at most \(\varepsilon\)
% can move the boundary of their convex hull by at most \(\varepsilon\).
% Since \(c(I)\) is the distance from \(p_G\) to the nearest boundary of
% this convex hull, its value can change by at most the same amount.

% To make this precise, fix a unit direction \(u\in\mathcal H\). The support
% function
% \[
% h_I(u)
% :=
% \max_{a\in I}\langle u,\mu_a\rangle
% \]
% describes how far the true convex hull extends in the direction \(u\).
% Similarly, define
% \[
% \widehat h_I(u)
% :=
% \max_{a\in I}\langle u,\widehat\mu_a\rangle.
% \]
% Since every vertex moves by at most \(\varepsilon\),
% \[
% |\langle u,\widehat\mu_a-\mu_a\rangle|
% \leq
% \varepsilon
% \]
% for every \(a\), and hence
% \[
% |\widehat h_I(u)-h_I(u)|
% \leq
% \varepsilon.
% \]

% For a fixed direction \(u\), the quantity
% \[
% h_I(u)-\langle u,p_G\rangle
% \]
% is the distance, in that direction, from \(p_G\) to the corresponding
% supporting boundary of the convex hull. The interior margin is obtained
% by taking the smallest such distance over all directions:
% \[
% c(I)
% =
% \min_{\substack{u\in\mathcal H\\\|u\|_2=1}}
% \left(
% h_I(u)-\langle u,p_G\rangle
% \right),
% \]
% and similarly
% \[
% \widehat c(I)
% =
% \min_{\substack{u\in\mathcal H\\\|u\|_2=1}}
% \left(
% \widehat h_I(u)-\langle u,p_G\rangle
% \right).
% \]

% Since the directional distance changes by at most \(\varepsilon\) for
% every \(u\), taking the minimum over \(u\) gives
% \[
% |\widehat c(I)-c(I)|
% \leq
% \varepsilon.
% \]
% \end{proof}

For a \(K\)-source support \(I\) such that
\[
p_G\in\operatorname{conv}\{\widehat\mu_a:a\in I\},
\]
define its empirical cost by
\[
\widehat B(I)
:=
\min\left\{
\sum_{a\in I}\ell_a q_a:
q_a\geq0,\quad
\sum_{a\in I}q_a=1,\quad
p_G=\sum_{a\in I}q_a\widehat\mu_a
\right\}.
\]

\begin{lemma}[Stability of the empirical cost]
\label{lem:cost-stability}
Let \(I\subseteq[m]\) be a \(K\)-source support with
\[
c(I)\geq c_0>0.
\]
Suppose that
\[
p_G\in
\operatorname{conv}\{\widehat\mu_a:a\in I\},
\]
so that the estimated source means can represent \(p_G\). If
\[
\max_{a\in I}
\|\widehat\mu_a-\mu_a\|_2
\leq
\varepsilon,
\]
then
\[
\left|
\widehat B(I)-B(I)
\right|
\leq
\frac{\ell_{\max}}{c_0}\varepsilon.
\]
\end{lemma}

\begin{proof}[Proof of Lemma~\ref{lem:cost-stability}] Let \(q\) be the unique probability vector satisfying
\[
p_G=\sum_{a\in I}q_a\mu_a,
\]
and let \(\widehat q\) be a probability vector attaining
\(\widehat B(I)\), so that
\[
p_G=\sum_{a\in I}\widehat q_a\widehat\mu_a.
\]

The argument is
\[
\widehat\mu_a\approx\mu_a
\Longrightarrow
\left\|
\sum_{a\in I}(\widehat q_a-q_a)\mu_a
\right\|_2
\leq
\varepsilon
\Longrightarrow
\widehat q\approx q
\Longrightarrow
\widehat B(I)\approx B(I),
\]
where the middle implication follows from
Lemma~\ref{lem:coefficient-stability}.\\

Note that since \(q\) and \(\widehat q\) are probability vectors,
\[
\sum_{a\in I}(\widehat q_a-q_a)=0.
\]

Moreover,
\[
\begin{aligned}
\sum_{a\in I}(\widehat q_a-q_a)\mu_a
&=
\sum_{a\in I}\widehat q_a\mu_a-p_G\\
&=
\sum_{a\in I}
\widehat q_a(\mu_a-\widehat\mu_a).
\end{aligned}
\]
Hence
\[
\left\|
\sum_{a\in I}(\widehat q_a-q_a)\mu_a
\right\|_2
\leq
\sum_{a\in I}
\widehat q_a\|\mu_a-\widehat\mu_a\|_2
\leq
\varepsilon.
\]

Applying Lemma~\ref{lem:coefficient-stability} with
\[
d_a=\widehat q_a-q_a
\]
gives
\[
\sum_{a\in I}
|\widehat q_a-q_a|
\leq
\frac{\varepsilon}{c_0}.
\]
Therefore
\[
\begin{aligned}
|\widehat B(I)-B(I)|
&=
\left|
\sum_{a\in I}\ell_a(\widehat q_a-q_a)
\right|\\
&\leq
\ell_{\max}
\sum_{a\in I}|\widehat q_a-q_a|\\
&\leq
\frac{\ell_{\max}}{c_0}\varepsilon.
\end{aligned}
\]
\end{proof}

\begin{lemma}[Coefficient stability from interiority]
\label{lem:coefficient-stability}
Let \(\mu_1,\ldots,\mu_K\in\mathbb R^d\), and define
\[
P:=\operatorname{conv}\{\mu_1,\ldots,\mu_K\}.
\]
Suppose \(\dim(P)=K-1\). Assume that there exist \(p\in P\) and
\(c>0\) such that
\[
\{x\in \operatorname{aff}(P):\|x-p\|_2\leq c\}
\subseteq P,
\]
where
\[
\operatorname{aff}(P)
=
\left\{
\sum_{i=1}^K \alpha_i \mu_i
:
\alpha_i\in\mathbb R,
\quad
\sum_{i=1}^K \alpha_i=1
\right\}.
\]

If \(d_1,\ldots,d_K\in\mathbb R\) satisfy
\[
\sum_{i=1}^K d_i=0,
\]
then
\[
\sum_{i=1}^K |d_i|
\leq
\frac{1}{c}
\left\|
\sum_{i=1}^K d_i\mu_i
\right\|_2.
\]
In particular, if
\[
\left\|
\sum_{i=1}^K d_i\mu_i
\right\|_2
\leq
\varepsilon,
\]
then
\[
\sum_{i=1}^K |d_i|
\leq
\frac{\varepsilon}{c}.
\]
\end{lemma}

\begin{proof}[Proof of Lemma~\ref{lem:coefficient-stability}]
Since \(\dim(P)=K-1\), the points
\(\mu_1,\ldots,\mu_K\) are affinely independent. Hence \(p\) has a
unique representation
\[
p=\sum_{i=1}^K q_i\mu_i,
\qquad
\sum_{i=1}^K q_i=1.
\]
Moreover, since a ball of positive radius centered at \(p\) is contained
in \(P\), the point \(p\) lies in the relative interior of \(P\), and therefore $q_i>0, \text{ for every } i.$
Set
\[
v:=\sum_{i=1}^K d_i\mu_i.
\]
For \(t\in\mathbb R\),
\[
p+t v
=
\sum_{i=1}^K(q_i+t d_i)\mu_i.
\]
Since
\[
\sum_{i=1}^K d_i=0,
\]
the coefficients satisfy
\[
\sum_{i=1}^K(q_i+t d_i)=1
\]
for every \(t\). Let
\[
M
:=
\max_{1\leq i\leq K}\frac{|d_i|}{q_i}.
\]
If \(M=0\), then all \(d_i=0\) and the result is immediate. Therefore \(M>0\), and choose \(i^\star\) attaining the maximum.

Choose \(t^\star\) with sign opposite to \(d_{i^\star}\) and magnitude
\[
|t^\star|
=
\frac{q_{i^\star}}{|d_{i^\star}|}
=
\frac{1}{M}.
\]
Then
\[
q_{i^\star}+t^\star d_{i^\star}=0.
\]
For every \(i\),
\[
|t^\star|
\leq
\frac{q_i}{|d_i|},
\]
with the convention that \(q_i/|d_i|=\infty\) when \(d_i=0\).
If \(t^\star d_i<0\), this implies
\[
q_i+t^\star d_i
\geq
q_i-|t^\star||d_i|
\geq0,
\]
while if \(t^\star d_i\geq0\), the coefficient is clearly nonnegative.
Thus all coefficients remain nonnegative, and at least one is zero. Therefore $p+t^\star v$
lies on the boundary of \(P\). Since the ball of radius \(c\) centered
at \(p\) is contained in \(P\), every boundary point of \(P\) is at
distance at least \(c\) from \(p\). Hence
\[
c
\leq
\|t^\star v\|_2
=
\frac{\|v\|_2}{M}.
\]
Thus
\[
M
\leq
\frac{\|v\|_2}{c}.
\]

By the definition of \(M\),
\[
|d_i|
\leq
\frac{\|v\|_2}{c}\,q_i
\]
for every \(i\). Summing and using \(\sum_i q_i=1\) gives
\[
\sum_{i=1}^K|d_i|
\leq
\frac{\|v\|_2}{c}
=
\frac{1}{c}
\left\|
\sum_{i=1}^K d_i\mu_i
\right\|_2.
\]

In particular, if
\[
\left\|
\sum_{i=1}^K d_i\mu_i
\right\|_2
\leq
\varepsilon,
\]
then
\[
\sum_{i=1}^K|d_i|
\leq
\frac{\varepsilon}{c}.
\]
\end{proof}

\begin{proof}[Proof of Theorem~\ref{thm:cost-aware}]
By Theorem~\ref{thm:cost-support}, with probability at least
\(1-\delta\), the selected support satisfies
\[
c(\widehat I)\ge \frac{c_0}{2},
\]
together with the stated bound on \(B(\widehat I)\).

Condition on the exploration history and on this event. After support
selection, Algorithm~\ref{alg:ucb-sampling} is initialized afresh and run
on the \(K\) sources in \(\widehat I\). In the categorical setting,
\(\|Y_t-p_G\|\le\sqrt2\). Applying the proof of
Theorem~\ref{thm:multidimensional-guarantee} with \(m=K\) and
\(c\ge c_0/2\) therefore gives
\[
\|\widehat p_T-p_G\|
=
O\!\left(
\frac{K^2}{c_0^2T}
\log\frac{K}{\delta c_0}
+
\frac{1}{c_0T}
\log\frac{1}{\eta c_0}
\right).
\]

It remains to control the realized cost. Let \(q\) be the unique mixture on
\(\widehat I\) satisfying
\[
p_G=\sum_{a\in\widehat I}q_a\mu_a,
\]
and let
\[
x_a:=\frac{N_a(T)}{T}
\]
be the realized sampling proportions during the post-exploration phase.
By Lemma~\ref{lem:coefficient-stability},
\[
\|x-q\|_1
\le
\frac{2}{c_0}
\left\|
\sum_{a\in\widehat I}x_a\mu_a-p_G
\right\|.
\]
Writing
\[
M_T:=\sum_{t=1}^T(Y_t-\mu_{A_t}),
\]
we have
\[
\sum_{a\in\widehat I}x_a\mu_a-p_G
=
\widehat p_T-p_G-\frac{M_T}{T}.
\]
Each coordinate of \(M_T\) is a martingale with bounded increments, so
\[
\frac{\|M_T\|}{T}
=
O_{\mathbb P}\!\left(\sqrt{\frac KT}\right).
\]
Hence
\[
\left|
\frac1T\sum_{t=1}^T\ell_{A_t}
-
B(\widehat I)
\right|
\le
\ell_{\max}\|x-q\|_1
=
O_{\mathbb P}\!\left(
\frac{\ell_{\max}}{c_0}\sqrt{\frac KT}
\right).
\]
Finally, the exploration phase costs at most \(mn\ell_{\max}\). Therefore
\[
\overline B_T
\le
B(\widehat I)
+
\frac{mn\ell_{\max}}{T}
+
O_{\mathbb P}\!\left(
\frac{\ell_{\max}}{c_0}\sqrt{\frac KT}
\right),
\]
as claimed.
\end{proof}
\begin{proof}[Proof of Proposition \ref{prop:one-dimensional-lower-bound}]
Condition on the history up to time \(T-1\), including the sampling
method \(A_T\) chosen for the final observation. Let
\[
N_{T-1}=\sum_{t=1}^{T-1}Y_t.
\]
Depending on whether \(Y_T=0\) or \(Y_T=1\), the final empirical
proportion is one of
\[
\frac{N_{T-1}}{T}
\qquad\text{and}\qquad
\frac{N_{T-1}+1}{T}.
\]
These two values are separated by \(1/T\). Hence they cannot both lie
within distance \(1/(2T)\) of \(p_G\). Therefore, for at least one of the
two possible values of \(Y_T\),
\[
\left|\widehat p_T-p_G\right|
\geq \frac{1}{2T}.
\]

Whichever sampling method is chosen at time \(T\), both outcomes
\(Y_T=0\) and \(Y_T=1\) have conditional probability at least \(q\).
Thus,
\[
\mathbb P\left(
\left|\widehat p_T-p_G\right|
\geq \frac{1}{2T}
\,\middle|\,\mathcal F_{T-1}
\right)
\geq q.
\]
Taking expectations gives the probability bound. The expectation bound
then follows from
\[
\mathbb E\left[
\left|\widehat p_T-p_G\right|
\right]
\geq
\frac{1}{2T}
\mathbb P\left(
\left|\widehat p_T-p_G\right|
\geq \frac{1}{2T}
\right).
\]
\end{proof}
\begin{proof}[Proof of Proposition \ref{prop:multidimensional-guarantee}]
Let
\[
p_G=\left(\frac12,\ldots,\frac12\right),
\]
and fix \(c\in(0,1/2)\). For every
\(\sigma\in\{-1,+1\}^K\), introduce a source with mean
\[
\mu_\sigma=p_G+c\sigma,
\]
for example by taking the coordinates of \(Y_t\) to be independent
Bernoulli random variables with
\[
\mathbb P(Y_{t,j}=1\mid A_t=\sigma)
=\frac12+c\sigma_j.
\]
The convex hull of the source means is the cube
\(p_G+[-c,c]^K\), and therefore contains the Euclidean ball
\(B(p_G,c)\).

Now condition on the history up to time \(T-1\) and on the source chosen
at time \(T\). Let
\[
N_{T-1,1}=\sum_{t=1}^{T-1}Y_{t,1}.
\]
According as \(Y_{T,1}=0\) or \(1\), the first coordinate of the final
empirical vector equals
\[
\frac{N_{T-1,1}}{T}
\qquad\text{or}\qquad
\frac{N_{T-1,1}+1}{T}.
\]
These two values are separated by \(1/T\), so at least one of them is at
distance at least \(1/(2T)\) from \(p_{G,1}=1/2\). Hence for at least one
of the two possible outcomes,
\[
\|\widehat p_T-p_G\|_2\geq \frac{1}{2T}.
\]

Under every source, both values \(Y_{T,1}=0\) and \(Y_{T,1}=1\) occur with
probability at least
\[
q:=\frac12-c.
\]
Therefore
\[
\mathbb P\left(
\|\widehat p_T-p_G\|_2\geq \frac{1}{2T}
\right)\geq q.
\]
The expectation bound follows immediately.
\end{proof}

\begin{lemma}[Cost-gap identification]
\label{lem:cost-gap}
Suppose \(I^\star\) and \(I\) satisfy
\[
B(I)
\geq
B(I^\star)+\Delta_B.
\]
Suppose further that
\[
|\widehat B(I)-B(I)|
\leq
\frac{\Delta_B}{4}
\]
and
\[
|\widehat B(I^\star)-B(I^\star)|
\leq
\frac{\Delta_B}{4}.
\]
Then
\[
\widehat B(I)
>
\widehat B(I^\star).
\]
\end{lemma}

\begin{proof}
We have
\[
\widehat B(I)
\geq
B(I)-\frac{\Delta_B}{4}
\]
and
\[
\widehat B(I^\star)
\leq
B(I^\star)+\frac{\Delta_B}{4}.
\]
Hence
\[
\begin{aligned}
\widehat B(I)-\widehat B(I^\star)
&\geq
B(I)-B(I^\star)
-\frac{\Delta_B}{2}\\
&\geq
\frac{\Delta_B}{2}
>0.
\end{aligned}
\]
\end{proof}

\begin{remark}
The support-selection error in Theorem~\ref{thm:cost-aware} decreases as
\(n^{-1/2}\).
In particular, for any \(\varepsilon>0\), it is enough to take
\[
n
\ge
\frac{9K\ell_{\max}^2}{2c_0^2\varepsilon^2}
\log\frac{2mK}{\delta}
\]
to guarantee
\[
B(\widehat I)
\le
\min_{\substack{|I|=K\\ c(I)\ge c_0}}
B(I)
+\varepsilon,
\]
in addition to the basic exploration requirement of the theorem.
This bound concerns the cost \(B(\widehat I)\) of the selected support;
the realized average cost \(\overline B_T\) additionally contains the
finite-\(T\) terms described in Theorem~\ref{thm:cost-aware}.

Let \(I^\star\) be the cheapest \(K\)-source support with
\(c(I^\star)\ge c_0\). If there exists \(\Delta_B>0\) such that
\[
B(I)\ge B(I^\star)+\Delta_B
\]
for every other \(K\)-source support \(I\) with
\(c(I)\ge c_0/2\), then exact recovery of \(I^\star\) is guaranteed whenever
\[
n
\ge
\frac{K}{c_0^2}
\log\frac{2mK}{\delta}
\max\left\{
8,\,
\frac{9\ell_{\max}^2}{2\Delta_B^2}
\right\}.
\]
\end{remark}

\section{A Martingale Central Limit Theorem for the Sampling Noise}
\label{app:martingale-clt}

This appendix studies the sampling-noise term
\[
M_T
:=
\sum_{t=1}^T
\bigl(Y_t-p_{A_t}\bigr)
\]
that appears in the representation-error decomposition. We prove that, under
a stabilization condition on the empirical sampling frequencies,
\[
\frac{M_T}{\sqrt{T}}
\]
converges in distribution to a Gaussian random variable.

\begin{lemma}[Sampling-noise scale under sign feedback]
\label{lem:sampling-noise-scale}
Assume
\[
p_P<p_G<p_O
\qquad\text{and}\qquad
p_P,p_O\in(0,1),
\]
and run Algorithm~\ref{alg:sign-based}. Then
\[
\frac{M_T}{\sqrt T}
\Longrightarrow
\mathcal N(0,\nu^\star),
\]
where
\[
\nu^\star
=
x^\star p_O(1-p_O)
+
(1-x^\star)p_P(1-p_P)
>0.
\]
Consequently, for every \(\varepsilon>0\), there exist constants
\(0<c_\varepsilon<C_\varepsilon<\infty\) such that, for all sufficiently
large \(T\),
\[
\mathbb P\left(
c_\varepsilon\sqrt T
\leq |M_T|
\leq C_\varepsilon\sqrt T
\right)
\geq 1-\varepsilon.
\]
\end{lemma}

\begin{proof}
Let
\[
\xi_t:=Y_t-p_{A_t},
\qquad
M_T=\sum_{t=1}^T\xi_t.
\]
If \(\mathcal F_{t-1}\) denotes the history before observing \(Y_t\), then
\(A_t\) is \(\mathcal F_{t-1}\)-measurable and
\[
\mathbb E[\xi_t\mid\mathcal F_{t-1}]
=
\mathbb E[Y_t-p_{A_t}\mid\mathcal F_{t-1}]
=
0.
\]
Thus \(M_T\) is a martingale with bounded increments, \(|\xi_t|\leq1\).
Moreover,
\[
\mathbb E[M_T^2]
=
\sum_{t=1}^T\mathbb E[\xi_t^2]
\leq \frac{T}{4},
\]
and hence
\[
\frac{M_T}{T}
\xrightarrow{\mathbb P}0.
\]

By Theorem~\ref{thm:one-dimensional-guarantee},
\[
S_T=O_{\mathbb P}(1),
\]
so \(S_T/T\to0\) in probability. The decomposition identity, multiplied by
\(T\), gives
\[
S_T
=
(p_O-p_P)\bigl(N_O(T)-x^\star T\bigr)
+
M_T.
\]
Dividing by \(T\) therefore yields
\[
\frac{N_O(T)}{T}
\xrightarrow{\mathbb P}
x^\star.
\]

The predictable quadratic variation of \(M_T\) is
\[
V_T
:=
\sum_{t=1}^T
\mathbb E[\xi_t^2\mid\mathcal F_{t-1}]
=
N_O(T)p_O(1-p_O)
+
N_P(T)p_P(1-p_P).
\]
Consequently,
\[
\frac{V_T}{T}
\xrightarrow{\mathbb P}
\nu^\star.
\]
Since the martingale increments are uniformly bounded, the Lindeberg
condition is automatic. The standard martingale central limit theorem then
gives
\[
\frac{M_T}{\sqrt T}
\Longrightarrow
\mathcal N(0,\nu^\star).
\]

Since \(\nu^\star>0\), the limiting Gaussian is nondegenerate. The final
claim follows by choosing \(c_\varepsilon\) sufficiently small and
\(C_\varepsilon\) sufficiently large and applying convergence in
distribution.
\end{proof}
\begin{proof}[Proof of Lemma~\ref{lem:representation-decomposition}]
Since
\[
\widehat p_T
=
\frac1T\sum_{t=1}^T Y_t,
\]
we write
\[
\widehat p_T-p_G
=
\frac1T\sum_{t=1}^T
\bigl(Y_t-p_{A_t}\bigr)
+
\frac1T\sum_{t=1}^T
\bigl(p_{A_t}-p_G\bigr).
\]
Now
\[
\frac1T\sum_{t=1}^T p_{A_t}
=
y_Tp_O+(1-y_T)p_P,
\]
while
\[
p_G=x^\star p_O+(1-x^\star)p_P.
\]
Therefore
\[
\frac1T\sum_{t=1}^T
\bigl(p_{A_t}-p_G\bigr)
=
(y_T-x^\star)(p_O-p_P),
\]
which gives the claimed identity.
\end{proof}

\section{Experimental Details and Additional Results}
\label{app:experiments}

This appendix documents the experiments of Section~\ref{sec:experiments}
and reports additional results. Every number is computed from the result
files written by the accompanying code, every run is seeded, and policies
compared within one experiment share the same seed.

\subsection{Data, Channels and Targets}
\label{app:exp-data}

\paragraph{Survey data.} We use the 2023 NHIS Sample Adult public-use file
\cite{nhis2023}, joined on the household identifier with the 2023 Paradata
file, which records the telephone status of each household (\(29{,}522\)
adults). The target population consists of the \(3{,}294\) adults who report
diagnosed diabetes (\texttt{DIBEV\_A}\(=1\)) and whose age group is known.
All compositions are weighted by the sample-adult weight
\texttt{WTFA\_A}. Fewer than \(300\) respondents report type~1 diabetes,
too few to estimate sixteen channel compositions, so the target
population is all adults with diagnosed diabetes.

\paragraph{Channels.} A channel is the set of target-population
respondents that a recruitment method can reach
(Table~\ref{tab:channels}); pools overlap, and pools with fewer than
\(150\) respondents would be discarded (none is). The source distribution of
a channel is the weighted distribution of its pool over the target
categories, so \(\mu_a\) is the weighted composition of the pool. Except
in Appendix~\ref{app:exp-downstream}, a recruit from channel \(a\) is a
simulated independent draw from \(\mu_a\), so \(T\) counts simulated
recruitment rounds and may exceed the number of NHIS respondents. In the
respondent-level experiment of Appendix~\ref{app:exp-downstream}, a draw
from a channel is an individual respondent, chosen with probability
proportional to its weight within the pool. Online health-information
seekers and patient-portal users cost \(\$72\) per enrolled participant and
all other channels \(\$199\): these are the medians per enrolled participant
for online and offline recruitment in the meta-analysis
of~\cite{brogger2020online}, whose ranges are \(\$4\)--\(\$251\) and
\(\$19\)--\(\$839\). The construction models coverage, i.e., whom a channel
can reach, and not differential willingness to enrol within a channel.

\begin{table}[!ht]
\caption{The sixteen recruitment channels (NHIS~2023 variable codes), their
number of respondents with diagnosed diabetes, their ethnic composition
(weighted, \%), and their cost per enrolled participant.}
\label{tab:channels}
\centering\small
\setlength{\tabcolsep}{3.5pt}
\begin{tabular}{llrrrrr}
\hline
channel & definition & \(n\) & Hispanic & Black & other & \(\$\)\\
\hline
online & looked up health information online, \texttt{HITLOOK\_A}\(=1\) & 1{,}583 & 13.6 & 13.9 & 72.4 & 72\\
portal & \texttt{HITCOMM\_A}\(=1\) or \texttt{HITTEST\_A}\(=1\) & 1{,}691 & 15.5 & 13.4 & 71.1 & 72\\
landline & dual, mostly or only landline, \texttt{AD\_PCLASS}\(\in\{3,4,5\}\) & 642 & 9.6 & 14.4 & 75.9 & 199\\
cell & only or mostly wireless, dual, \texttt{AD\_PCLASS}\(\in\{1,2,3\}\) & 2{,}746 & 20.0 & 15.1 & 64.8 & 199\\
Medicare & \texttt{MEDICARE\_A}\(\in\{1,2\}\) & 2{,}020 & 14.3 & 14.6 & 71.1 & 199\\
Medicaid & \texttt{MEDICAID\_A}\(\in\{1,2\}\) & 637 & 30.8 & 20.7 & 48.6 & 199\\
VA & \texttt{VAHOSP\_A}\(=1\) & 198 & 15.5 & 19.5 & 65.0 & 199\\
emergency & emergency visit in past 12 months, \texttt{EMERG12MTC\_A}\(\in\{1,\dots,4\}\) & 1{,}021 & 19.8 & 18.1 & 62.0 & 199\\
inpatient & overnight hospital stay, \texttt{HOSPONGT\_A}\(=1\) & 619 & 17.1 & 15.6 & 67.3 & 199\\
telehealth & \texttt{VIRAPP12M\_A}\(=1\) & 1{,}123 & 18.4 & 16.5 & 65.1 & 199\\
Spanish-language & \texttt{LANGSPECR\_A}\(=1\) & 289 & 94.2 & 0.4 & 5.4 & 199\\
clinics, Northeast & \texttt{USUALPL\_A}\(=1\), \texttt{LASTDR\_A}\(=1\), \texttt{REGION}\(=1\) & 435 & 19.2 & 12.5 & 68.3 & 199\\
clinics, Midwest & same, \texttt{REGION}\(=2\) & 701 & 4.3 & 10.7 & 85.0 & 199\\
clinics, South & same, \texttt{REGION}\(=3\) & 1{,}253 & 15.9 & 24.5 & 59.6 & 199\\
clinics, West & same, \texttt{REGION}\(=4\) & 691 & 37.2 & 4.8 & 57.9 & 199\\
rural & nonmetropolitan, \texttt{URBRRL}\(=4\) & 666 & 5.2 & 9.2 & 85.6 & 199\\
\hline
target & diagnosed diabetes, \texttt{DIBEV\_A}\(=1\) & 3{,}294 & 19.0 & 15.5 & 65.4 & \\
\hline
\end{tabular}
\end{table}

\paragraph{Targets.} Table~\ref{tab:targets} lists the targets. The two
two-source instances use a binary label: the share aged \(65+\) with
channels online (\(40.0\%\)) and landline (\(78.0\%\)), target \(47.1\%\);
and the non-Hispanic Black share with channels online (\(13.9\%\)) and
clinics in the South (\(24.5\%\)), target \(15.5\%\). All other targets use
the sixteen channels. For the marginal goals, the label
\(Y_t\in\{0,1\}^4\) indicates whether a recruit is Black, Hispanic, aged
\(65+\) and female, so \(\|Y_t\|_0\le\ell_0=3\), and the margin is computed
in \(\mathbb R^4\); for categorical targets it is computed within the
affine hull of the simplex, as in Section~\ref{sec:cost}.

\begin{table}[!ht]
\caption{Targets: categories, target composition \(p_G\) (\%), and margin
\(c\) of \(p_G\) in the convex hull of the channel means.}
\label{tab:targets}
\centering\small
\begin{tabular}{llrr}
\hline
target & categories and \(p_G\) (\%) & \(K\) & \(c\)\\
\hline
two sources, age & \(65+\): 47.1, \(<65\): 52.9 (online, landline) & 2 & 0.1001\\
two sources, Black & Black: 15.5, other: 84.5 (online, clinics South) & 2 & 0.0228\\
\hline
Black / other & Black: 15.5, other: 84.5 & 2 & 0.1275\\
ethnicity & Hispanic: 19.0, non-Hispanic Black: 15.5, other: 65.4 & 3 & 0.0822\\
age & 18--49: 20.0, 50--64: 32.9, \(65+\): 47.1 & 3 & 0.0469\\
age\(\times\)sex & men \(<65\): 26.5, men \(65+\): 23.5, women \(<65\): 26.4, women \(65+\): 23.6 & 4 & 0.0256\\
Black\(\times\)age & Black \(<65\): 8.7, Black \(65+\): 6.8, other \(<65\): 44.3, other \(65+\): 40.2 & 4 & 0.0086\\
ethnicity\(\times\)age & H \(<65\): 12.1, H \(65+\): 7.0, B \(<65\): 8.7, B \(65+\): 6.8, O \(<65\): 32.2, O \(65+\): 33.3 & 6 & 0.0014\\
marginal goals & Black: 15.5, Hispanic: 19.0, \(65+\): 47.1, female: 50.0 (multi-label) & 4 & 0.0537\\
\hline
\end{tabular}
\end{table}

\FloatBarrier
\subsection{Policies and Protocol}
\label{app:exp-impl}

\paragraph{Algorithms.} Algorithm~\ref{alg:sign-based} is run as stated,
with \(P\) the channel whose share of the designated group is below
\(p_G\). Algorithm~\ref{alg:ucb-sampling} is run as stated, with
\(\delta=0.05\) and the radius \(r_a(n)\) of Appendix~A; unsampled
channels have index \(-\infty\) (ties are broken by channel order), and
when \(S_t=0\) the least-sampled channel is chosen. Its variants multiply
\(r_a\) by \(\beta\in\{0.3,0.1\}\) or remove it (\(\beta=0\), the greedy
plug-in rule). A value \(\beta<1\) is not a re-tuned confidence level:
since \(r_a(n)\propto\sqrt{\log(A/\delta)}\) with \(A=\pi^2mKn^2/3\),
matching \(\beta r_a(n)\) would require \(\delta'=A^{1-\beta^2}\delta^{\beta^2}\),
e.g., \(\delta'\approx2\times10^7\) for \(\beta=0.3\), \(K=3\), \(m=16\),
\(n=10^3\). The variants therefore have no guarantee from
Theorem~\ref{thm:multidimensional-guarantee}, and we report them as
empirical results only. A round is counted as uncertain when
\(r_{A_{t+1}}(N_{A_{t+1}}(t))>c/4\) for the unscaled radius.
Algorithm~\ref{alg:cost-support} samples every channel \(n\) times in
rotation, and these samples are kept in the final sample. For every
\(K\)-subset \(I\), it computes the barycentric coordinates \(q(I)\) of
\(p_G\) with respect to the estimated means and the margin
\(\widehat c(I)=\min_{i\in I}q_i(I)h_i(I)\), where \(h_i(I)\) is the distance
from \(\widehat\mu_i\) to the affine hull of the other \(K-1\) estimated
means, so that \(q_i(I)h_i(I)\) is the distance from \(p_G\) to the facet
opposite \(\widehat\mu_i\). A support is admissible when \(q(I)\ge0\) and
\(\widehat c(I)\ge3c_0/4\). Because \(K\) affinely independent means
determine \(q(I)\), the linear program of Algorithm~\ref{alg:cost-support}
reduces to evaluating \(\sum_{a\in I}\ell_aq_a(I)\). If Algorithm~\ref{alg:cost-support} returns
\textsc{Fail}, we run Algorithm~\ref{alg:ucb-sampling} on all channels. After selection, Algorithm~\ref{alg:ucb-sampling} runs on
\(\widehat I\) with \(m=K\) in the radius and \(\beta=0.3\). The margin
\(c\) of a target in the hull of all channels is the smallest distance from
\(p_G\) to a facet of that hull.

\paragraph{Baselines.} The oracle closed loop chooses
\(\arg\min_a\langle S_t,\mu_a-p_G\rangle\) with the true means. The
oracle open loop follows a fixed mixture \(x^\star\) with
\(\sum_ax^\star_a\mu_a=p_G\) (unique for two sources; the cheapest one
otherwise) by deterministic herding,
\(A_{t+1}\in\arg\max_a\{(t+1)x^\star_a-N_a(t)\}\), so that the share of
recruits from each channel is within \(O(1/T)\) of \(x^\star_a\); the
allocation is fixed in advance and does not react to the realized
compositions; its i.i.d.\ variant
draws \(A_t\sim x^\star\) independently. Explore-then-commit samples every
channel \(\lfloor0.1\,T/m\rfloor\) times in rotation (kept in the sample),
then computes the cheapest mixture that makes the expected final
composition equal to \(p_G\) under the estimated means and the labels
already collected (or, if none exists, the mixture closest to it in
\(\ell_\infty\)), and follows it by herding. Because it is configured for
a horizon \(T\), we evaluate it only at that horizon. The myopic cost-ratio
rule of Section~\ref{sec:cost} chooses
\(\arg\min_a\langle S_t,\mu_a-p_G\rangle/\ell_a\). In the NHIS cost
experiments it does not know the means, like our algorithms: it replaces
\(\mu_a\) by \(\widehat\mu_a\) and subtracts the exploration bonus of
Algorithm~\ref{alg:ucb-sampling} (radius \(0.3\,r_a\)), both divided by
\(\ell_a\), and samples every channel once before using the index. The
simulation of Example~\ref{ex:myopic} also runs it with the true means. No
policy discards recruits. The naive designs sample one fixed channel, rotate over 
all channels, or follow the mixture proportional to the weighted pool
sizes.

\paragraph{Protocol and metrics.} Table~\ref{tab:protocol} lists
horizons and replications. The error is \(\|\widehat p_T-p_G\|_2\),
averaged over replications, and \(\|S_T\|=T\|\widehat p_T-p_G\|\). For
two sources we report \(|\widehat p_T-p_G|\) for the designated group,
which equals \(\|\widehat p_T-p_G\|_2/\sqrt2\). Errors are recorded on a
logarithmic grid of horizons. Slopes are least-squares fits of the log mean
error against \(\log T\) for \(T\ge10^3\) (two sources) or
\(T\ge3\times10^3\) (many sources). \(N_{90}\) is the first grid point
(26 points from \(10\) to \(10^5\)) at which the \(90\)th percentile of the
error is at most \(0.01\).

\begin{table}[!ht]
\caption{Experiments, horizons and replications.}
\label{tab:protocol}
\centering\small
\setlength{\tabcolsep}{3.5pt}
\begin{tabular}{lllrr}
\hline
experiment & target & policies & \(T\) & runs\\
\hline
two sources (Fig.~\ref{fig:main}(a), \ref{fig:app-two}) & age; Black & Alg.~\ref{alg:sign-based}, oracle open loop, ETC, naive & \(10^5\) & 2{,}000\\
robustness (Fig.~\ref{fig:app-robust}) & age; Black & Alg.~\ref{alg:sign-based} with modifications, oracle open loop, ETC & \(2\times10^4\) & 500\\
many sources (Fig.~\ref{fig:main}(b), Table~\ref{tab:many}) & ethnicity & Alg.~\ref{alg:ucb-sampling} (\(\beta\in\{1,0.3,0.1,0\}\)), oracles, ETC, naive & \(10^5\) & 250\\
marginal goals (Fig.~\ref{fig:app-many}(a)) & marginal goals & as above & \(10^5\) & 250\\
Black\(\times\)age (Fig.~\ref{fig:app-many}(b)) & Black\(\times\)age & Alg.~\ref{alg:ucb-sampling} (\(\beta=0.3\)), oracle closed loop, naive & \(10^5\) & 250\\
cost per recruit (Fig.~\ref{fig:main}(c), \ref{fig:app-budget}(a,b), Table~\ref{tab:cost}) & Black/other & Table~\ref{tab:cost} & \(3\times10^5\) & 200\\
cost, fixed budget (Fig.~\ref{fig:app-budget}(c)) & Black/other & Alg.~\ref{alg:cost-support}\(\to\)\ref{alg:ucb-sampling}, myopic rule, Alg.~\ref{alg:ucb-sampling}, oracle & \(\$1\)--\(16\)M & 200\\
cost, \(c_0\) sweep (Table~\ref{tab:c0}) & Black/other & Alg.~\ref{alg:cost-support}\(\to\)\ref{alg:ucb-sampling}, \(n\in\{400,800,1{,}600\}\) & \(10^5\) & 100\\
support selection (Table~\ref{tab:selection}) & Black/other & Alg.~\ref{alg:cost-support}, \(n\in\{25,\dots,1{,}600\}\), 10 cost vectors & -- & 250\\
Example~\ref{ex:myopic} & synthetic & myopic rule (true and estimated means), Alg.~\ref{alg:sign-based}, Alg.~\ref{alg:cost-support}\(\to\)\ref{alg:ucb-sampling} & \(10^5\) & 100\\
downstream (Table~\ref{tab:downstream}) & ethnicity & Alg.~\ref{alg:ucb-sampling} (\(\beta=0.3\)), oracle closed loop, naive & \(2\times10^4\) & 200\\
\hline
\end{tabular}
\end{table}

\FloatBarrier
\subsection{Two Sources}
\label{app:exp-two}

Figure~\ref{fig:app-two} and Table~\ref{tab:two} show both two-source
instances. On both, Algorithm~\ref{alg:sign-based} has slope \(-1\) and a
bounded imbalance, \(E|S_T|=1.43\) and \(3.57\) at \(T=10^5\). These values
are \(13\) and \(51\) times the lower bound \(q/2\) of
Proposition~\ref{prop:one-dimensional-lower-bound} (\(0.110\) and \(0.070\))
and \(10\) and \(11\) times smaller than the value \((\log C+1)/\lambda\)
implied by Theorem~\ref{thm:one-dimensional-guarantee}, with \(\lambda\) and
\(C\) from its proof and \(\lambda\) optimized (\(14.7\) and \(40.8\)). The
oracle open loop matches the prediction of
Lemma~\ref{lem:sampling-noise-scale}: \(\sqrt T\,E|\widehat p_T-p_G|\) at
\(T=10^5\) is \(0.378\) and \(0.279\), against \(\sqrt{2\nu^\star/\pi}=0.380\)
and \(0.287\). Deterministic herding is only marginally better than drawing channels
independently from \(x^\star\) (errors \(1.2\) and \(1.3\times10^{-3}\) at
\(T=10^5\) for the age instance): removing the allocation error hardly
helps when the sampling error is of order \(T^{-1/2}\). At \(T=10^5\), the
allocation and sampling errors of Lemma~\ref{lem:representation-decomposition}
have correlation \(-0.99993\) and \(-0.9993\) across replications under
Algorithm~\ref{alg:sign-based}, and \(-0.25\) and \(-0.28\) under ETC.

\begin{figure}[!ht]
\centering
\includegraphics[width=\textwidth]{figA_twosource}
\caption{Two sources. Left: error; right: imbalance \(E|S_T|\), with the
lower bound \(q/2\) of Proposition~\ref{prop:one-dimensional-lower-bound}
(dotted) and the value implied by Theorem~\ref{thm:one-dimensional-guarantee}
(dashed). The thin green line \(\sqrt{2\nu^\star/(\pi T)}\) lies on top of
the oracle open-loop curves.}
\label{fig:app-two}
\end{figure}

\begin{table}[!ht]
\caption{Two sources, 2{,}000 replications: error at \(T=853\) and
\(T=10^5\), imbalance \(E|S_T|\) at \(T=10^5\), and slope of the error for
\(T\ge10^3\). ETC is evaluated only at its horizon \(10^5\). ``Offline
only'' is landline (age) or clinics in the South (Black).}
\label{tab:two}
\centering\small
\begin{tabular}{l|rrrr|rrrr}
\hline
 & \multicolumn{4}{c|}{age: online vs landline} & \multicolumn{4}{c}{Black: online vs clinics, South}\\
policy & \(T=853\) & \(T=10^5\) & \(E|S_T|\) & slope & \(T=853\) & \(T=10^5\) & \(E|S_T|\) & slope\\
\hline
Algorithm~\ref{alg:sign-based} & 0.0018 & \(1.4\cdot10^{-5}\) & 1.43 & \(-1.01\) & 0.0038 & \(3.6\cdot10^{-5}\) & 3.57 & \(-0.99\)\\
oracle open loop & 0.0132 & \(1.2\cdot10^{-3}\) & 119 & \(-0.50\) & 0.0099 & \(8.8\cdot10^{-4}\) & 88 & \(-0.51\)\\
oracle open loop, i.i.d. & 0.0137 & \(1.3\cdot10^{-3}\) & 126 & \(-0.50\) & 0.0099 & \(9.1\cdot10^{-4}\) & 91 & \(-0.50\)\\
ETC & -- & \(4.5\cdot10^{-3}\) & 447 & -- & -- & \(3.2\cdot10^{-3}\) & 325 & --\\
online only & 0.0703 & \(7.1\cdot10^{-2}\) & 7{,}081 & \(0.00\) & 0.0174 & \(1.6\cdot10^{-2}\) & 1{,}609 & \(0.00\)\\
offline only & 0.3097 & \(3.1\cdot10^{-1}\) & 30{,}921 & \(0.00\) & 0.0902 & \(9.0\cdot10^{-2}\) & 9{,}018 & \(0.00\)\\
uniform & 0.1194 & \(1.2\cdot10^{-1}\) & 11{,}920 & \(0.00\) & 0.0369 & \(3.7\cdot10^{-2}\) & 3{,}704 & \(0.00\)\\
\hline
\end{tabular}
\end{table}

\paragraph{Tail of the imbalance.} Figure~\ref{fig:app-tails} shows the
empirical tail \(P(|S_T|\ge k)\) under Algorithm~\ref{alg:sign-based}. It
does not change with \(T\) for \(T\ge10^3\) and decays geometrically in
\(k\), as Theorem~\ref{thm:one-dimensional-guarantee} states. The bound
\(\inf_\lambda C(\lambda)e^{-\lambda k}\), with \(C(\lambda)\) from the
proof, is valid but loose: it equals one for \(k\le12\) (age) and for all
\(k\le24\) (Black), whereas the empirical \(P(|S_T|\ge10)\) at \(T=10^5\) is
\(0.0035\) and \(0.067\).

\begin{figure}[!ht]
\centering
\includegraphics[width=\textwidth]{figA_tails}
\caption{Empirical tail of \(|S_T|\) under Algorithm~\ref{alg:sign-based}
(2{,}000 replications, so probabilities below \(5\times10^{-4}\) are not
resolved) and the bound of Theorem~\ref{thm:one-dimensional-guarantee}.}
\label{fig:app-tails}
\end{figure}

\paragraph{Robustness.} Figure~\ref{fig:app-robust} and
Table~\ref{tab:robust} stress Algorithm~\ref{alg:sign-based} on the age
instance. (a)~When the ordering assumption fails, with both channels below
the target (patient portal \(13.4\%\) and online \(13.9\%\) Black, target
\(15.5\%\)), the error converges to \(1.6\) pp, the distance from the target
to the closer channel, and does not diverge. (b)~Moving both channel means
toward the target by a factor \(s\) multiplies \(E|S_T|\) by about \(1/s\)
(\(1.4\), \(3.0\), \(6.2\) and \(12.3\) for \(s=1,\frac12,\frac14,\frac18\)).
(c)~If decisions are revised only every \(B\) recruits, or labels arrive
\(d\) rounds late, the excess of \(E|S_T|\) over the reference grows
roughly linearly in \(B\) and \(d\) (by \(1.1\) and \(6.2\) for
\(B=10,50\); by \(1.3\) and \(6.2\) for \(d=10,50\)), while the slope stays
close to \(-1\). (d)~When the \(65+\) shares of both
channels drift linearly by 5 pp over the horizon (online from \(40\%\) to
\(45\%\), landline from \(78\%\) to \(73\%\)), Algorithm~\ref{alg:sign-based}
ends with error \(2.6\times10^{-4}\). Its \(E|S_T|\) grows from about
\(1.5\) to \(5.3\) because the drift shrinks the gap between the lower
channel and the target from \(0.071\) to \(0.021\), as in~(b). The oracle
open loop and ETC fix their allocations for the initial means and end
\(1.6\) and \(1.7\) pp from the target.

\begin{figure}[!ht]
\centering
\includegraphics[width=\textwidth]{figA_robust}
\caption{Algorithm~\ref{alg:sign-based} under violated assumptions and
operational constraints (500 replications, \(T=2\times10^4\)). In (d), ETC
is evaluated at its horizon only.}
\label{fig:app-robust}
\end{figure}

\begin{table}[!ht]
\caption{Robustness of Algorithm~\ref{alg:sign-based} at
\(T=2\times10^4\) (500 replications): imbalance, error, and slope of the
error for \(T\ge10^3\).}
\label{tab:robust}
\centering\small
\begin{tabular}{lrrr}
\hline
variant & \(E|S_T|\) & \(|\widehat p_T-p_G|\) & slope\\
\hline
reference & 1.40 & \(7.0\cdot10^{-5}\) & \(-1.03\)\\
separation \(\times\frac12\) & 3.03 & \(1.5\cdot10^{-4}\) & \(-0.98\)\\
separation \(\times\frac14\) & 6.19 & \(3.1\cdot10^{-4}\) & \(-0.97\)\\
separation \(\times\frac18\) & 12.30 & \(6.1\cdot10^{-4}\) & \(-0.86\)\\
decisions in batches of 10 & 2.53 & \(1.3\cdot10^{-4}\) & \(-0.98\)\\
decisions in batches of 50 & 7.59 & \(3.8\cdot10^{-4}\) & \(-0.97\)\\
labels delayed by 10 & 2.69 & \(1.3\cdot10^{-4}\) & \(-1.00\)\\
labels delayed by 50 & 7.61 & \(3.8\cdot10^{-4}\) & \(-0.99\)\\
drift: Algorithm~\ref{alg:sign-based} & 5.26 & \(2.6\cdot10^{-4}\) & \(-0.62\)\\
drift: oracle open loop & 314 & \(1.6\cdot10^{-2}\) & \(0.12\)\\
drift: ETC & 340 & \(1.7\cdot10^{-2}\) & --\\
ordering violated & 322 & \(1.6\cdot10^{-2}\) & \(-0.03\)\\
\hline
\end{tabular}
\end{table}

\FloatBarrier
\subsection{Many Sources}
\label{app:exp-many}

\paragraph{Ethnic target.} Table~\ref{tab:many} gives the numbers behind
Figure~\ref{fig:main}(b). Without an exploration bonus (\(r_a\equiv0\)), the
algorithm trusts its current estimates: a channel that looks unhelpful after
its first few samples may be sampled too little afterward to correct the
estimate, and at least \(10\%\) of runs finish \(0.03\) or more from
\(p_G\). Some exploration is therefore necessary, but the full theoretical
radius is larger than needed here: radius \(0.3\,r_a\) reaches the
\(N_{90}\) criterion three times sooner.

\begin{table}[!ht]
\caption{Ethnic target, 16 channels (250 replications). Error
\(\|\widehat p_T-p_G\|_2\) at \(T=832\) (ETC still explores); \(N_{90}\):
first logged \(T\) with error \(\le0.01\) in \(90\%\) of runs;
\(\|S_T\|_2\) and cost per recruit at \(T=10^5\).}
\label{tab:many}
\centering\small
\begin{tabular}{lrrrr}
\hline
policy & error & \(N_{90}\) & \(\|S_T\|_2\) & \(\$\)/recruit\\
\hline
Algorithm~\ref{alg:ucb-sampling} & 0.0116 & 3{,}631 & 2.5 & 193\\
\quad radius \(0.3\,r_a\) & 0.0071 & 1{,}202 & 1.9 & 198\\
\quad greedy, \(r_a\equiv0\) & 0.0125 & \(>10^5\) & 532 & 192\\
oracle closed loop & 0.0024 & 398 & 2.0 & 199\\
oracle open loop & 0.0211 & 15{,}849 & 194 & 106\\
ETC & -- & \(>10^5\) & 1{,}731 & 108\\
online only & 0.0907 & \(>10^5\) & 8{,}988 & 72\\
uniform & 0.0402 & \(>10^5\) & 3{,}517 & 183\\
\hline
\end{tabular}
\end{table}

\paragraph{Marginal goals.} Diversity goals are often stated per
attribute rather than per cell. With the multi-label target of
Table~\ref{tab:targets} (\(c=0.054\)), the picture of
Section~\ref{sec:experiments} is unchanged (Figure~\ref{fig:app-many}(a),
Table~\ref{tab:many-extra}). Algorithm~\ref{alg:ucb-sampling} and its
variants with a smaller radius end within \(15\%\) of the oracle closed
loop at \(T=10^5\) (\(\|S_T\|\) of \(4.4\) to \(4.8\) against \(4.2\)). The
greedy rule has mean \(\|S_T\|=110\), driven by a minority of stalled runs
(its \(90\)th percentile is \(11.6\)). The oracle open loop reaches
\(\|S_T\|=246\).

\paragraph{Black\(\times\)age.} With \(c=0.0086\), the radius-\(0.3\)
variant of Algorithm~\ref{alg:ucb-sampling} is still in its transient at
\(T=10^5\) (\(\|S_T\|=35\) against \(7.8\) for the oracle), although its
error at \(T=832\), \(0.014\), is a third of that of uniform recruitment
(Figure~\ref{fig:app-many}(b), Table~\ref{tab:many-extra}).

\begin{figure}[!ht]
\centering
\includegraphics[width=\textwidth]{figA_multidim}
\caption{Many sources. (a) Marginal goals (multi-label, \(\ell_0=3\), 250
replications). (b) Black\(\times\)age (250 replications). (c) Error over
\(T\) for the six targets of Figure~\ref{fig:main}(c) (120 replications);
solid: Algorithm~\ref{alg:ucb-sampling} with radius \(0.3\,r_a\), dotted:
oracle closed loop.}
\label{fig:app-many}
\end{figure}

\begin{table}[!ht]
\caption{Marginal goals and Black\(\times\)age: error at \(T=832\), mean and
\(90\)th percentile of \(\|S_T\|\) at \(T=10^5\), and cost per recruit at
\(T=10^5\).}
\label{tab:many-extra}
\centering\small
\begin{tabular}{llrrrr}
\hline
target & policy & error & \(\|S_T\|\) & \(90\%\) & \(\$\)/recruit\\
\hline
marginal goals & Algorithm~\ref{alg:ucb-sampling} & 0.0233 & 4.8 & 8.6 & 179\\
 & \quad radius \(0.3\,r_a\) & 0.0117 & 4.4 & 7.7 & 181\\
 & \quad radius \(0.1\,r_a\) & 0.0101 & 4.5 & 8.1 & 184\\
 & \quad greedy, \(r_a\equiv0\) & 0.0121 & 110.1 & 11.6 & 184\\
 & oracle closed loop & 0.0049 & 4.2 & 7.4 & 186\\
 & oracle open loop & 0.0278 & 246.5 & 358.5 & 124\\
 & ETC & -- & 2{,}005.1 & 2{,}939.2 & 126\\
 & online only & 0.0966 & 9{,}238.7 & 9{,}403.5 & 72\\
 & uniform & 0.0515 & 4{,}402.0 & 4{,}581.2 & 183\\
 & pool-proportional & 0.0375 & 2{,}598.7 & 2{,}784.6 & 173\\
\hline
Black\(\times\)age & Alg.~\ref{alg:ucb-sampling}, radius \(0.3\,r_a\) & 0.0140 & 35.0 & 81.0 & 190\\
 & oracle closed loop & 0.0076 & 7.8 & 15.9 & 197\\
 & online only & 0.0825 & 8{,}063.9 & 8{,}331.8 & 72\\
 & uniform & 0.0438 & 3{,}963.3 & 4{,}213.4 & 183\\
\hline
\end{tabular}
\end{table}

\paragraph{The margin ladder.} Figure~\ref{fig:app-ladder} and
Table~\ref{tab:ladder} show the imbalance at \(T=10^5\) against the margin
for the six targets, from Black/other (\(K=2\), \(c=0.128\)) to
ethnicity\(\times\)age (\(K=6\), \(c=0.0014\)). The product \(c\|S_T\|\)
lies in \([0.07,0.25]\) for the oracle closed loop and in \([0.12,0.42]\)
for Algorithm~\ref{alg:ucb-sampling} with radius \(0.3\,r_a\); it varies by
a factor of at most \(3.6\) while \(c\) varies by a factor of \(92\). Finer
targets are also harder to match early on: after \(785\) recruits,
Algorithm~\ref{alg:ucb-sampling} is \(5.8\) times closer to the target than
uniform recruitment for ethnicity, but only \(2.7\) times closer for
ethnicity\(\times\)age. The gap between
Algorithm~\ref{alg:ucb-sampling} and the oracle widens as \(c\) decreases,
from none at \(c=0.128\) to a factor of \(4\) to \(5\) for \(c<0.01\),
where the learning transient has not ended by \(T=10^5\)
(Figure~\ref{fig:app-many}(c)).

\begin{figure}[!ht]
\centering
\includegraphics[width=0.5\textwidth]{figA_ladder}
\caption{Imbalance \(\|S_T\|_2\) at \(T=10^5\) against the margin \(c\) for
six targets over the same sixteen channels (right to left: Black/other,
ethnicity, age, age\(\times\)sex, Black\(\times\)age, ethnicity\(\times\)age;
120 runs). Algorithm~\ref{alg:ucb-sampling} uses radius \(0.3\,r_a\).}
\label{fig:app-ladder}
\end{figure}

\begin{table}[!ht]
\caption{Margin ladder over the sixteen channels (120 replications):
error at \(T=785\), \(\|S_T\|\) and \(c\|S_T\|\) at \(T=10^5\).
Algorithm~\ref{alg:ucb-sampling} uses radius \(0.3\,r_a\).}
\label{tab:ladder}
\centering\small
\begin{tabular}{lrr|rrr|rr|rr}
\hline
 & & & \multicolumn{3}{c|}{error at \(T=785\)} & \multicolumn{2}{c|}{\(\|S_T\|\)} & \multicolumn{2}{c}{\(c\|S_T\|\)}\\
target & \(K\) & \(c\) & Alg.~\ref{alg:ucb-sampling} & oracle & uniform & Alg.~\ref{alg:ucb-sampling} & oracle & Alg.~\ref{alg:ucb-sampling} & oracle\\
\hline
Black / other & 2 & 0.1275 & 0.0037 & 0.0011 & 0.0220 & 0.9 & 0.9 & 0.12 & 0.12\\
ethnicity & 3 & 0.0822 & 0.0070 & 0.0025 & 0.0407 & 1.9 & 1.8 & 0.15 & 0.15\\
age & 3 & 0.0469 & 0.0115 & 0.0042 & 0.0412 & 3.4 & 3.1 & 0.16 & 0.15\\
age\(\times\)sex & 4 & 0.0256 & 0.0240 & 0.0099 & 0.0439 & 16.3 & 9.9 & 0.42 & 0.25\\
Black\(\times\)age & 4 & 0.0086 & 0.0137 & 0.0073 & 0.0454 & 41.4 & 8.4 & 0.36 & 0.07\\
ethnicity\(\times\)age & 6 & 0.0014 & 0.0206 & 0.0113 & 0.0563 & 261.4 & 60.8 & 0.36 & 0.08\\
\hline
\end{tabular}
\end{table}

\paragraph{Constants of Theorem~\ref{thm:multidimensional-guarantee}.}
Table~\ref{tab:constants} evaluates the constants of
Remark~\ref{rem:constants} with \(\delta=\eta=0.05\) and with
\(R^2=\max_y\|y-p_G\|^2\) over the possible labels \(y\) (\(1.18\), \(1.23\)
and \(1.28\)), which is smaller than \(\ell_0+1\), so the table
understates the bound as stated. In every case \(N(c)\) exceeds the
horizon, so every round counts as uncertain in the proof. Even
\(r_a(10^5)=0.022\) exceeds \(c/4=0.021\) for the ethnic target. The bound
exceeds the measured imbalance by factors of \(8\times10^5\) to
\(9\times10^6\).

\begin{table}[!ht]
\caption{Constants of Theorem~\ref{thm:multidimensional-guarantee}
(Remark~\ref{rem:constants}) and measured imbalance at \(T=10^5\).}
\label{tab:constants}
\centering\small
\begin{tabular}{lrrrrrr}
\hline
target & \(c\) & \(N(c)\) & \(M(c)\) & bound on \(\|S_T\|\) & measured \(\|S_T\|\) & policy measured\\
\hline
ethnicity & 0.082 & \(1.1\cdot10^5\) & \(1.8\cdot10^6\) & \(2.0\cdot10^6\) & 2.5 & Algorithm~\ref{alg:ucb-sampling}\\
Black\(\times\)age & 0.0086 & \(1.8\cdot10^7\) & \(2.9\cdot10^8\) & \(3.2\cdot10^8\) & 35.0 & radius \(0.3\,r_a\)\\
marginal goals & 0.054 & \(3.9\cdot10^5\) & \(6.2\cdot10^6\) & \(7.0\cdot10^6\) & 4.8 & Algorithm~\ref{alg:ucb-sampling}\\
\hline
\end{tabular}
\end{table}

\FloatBarrier
\subsection{Cost-Aware Sampling}
\label{app:exp-cost}

All cost-aware experiments use the Black/other target (Black: \(15.5\%\),
other: \(84.5\%\); \(K=2\), margin \(c=0.128\) over all channels) and the
median costs of the meta-analysis, \(\$72\) for online and \(\$199\) for
offline channels. No policy discards recruits. The myopic rule uses
estimated means, and Algorithm~\ref{alg:ucb-sampling} after
Algorithm~\ref{alg:cost-support} uses the exploration samples, which stay in
the sample, to estimate the means.

\paragraph{Supports.} For \(c_0=0.02\), \(28\) two-channel supports have
margin at least \(c_0\). The cheapest is \(I^\star=\{\)online, clinics
South\(\}\), with weights \((0.85,0.15)\), margin \(0.023\) and cost
\(\$91.27\); it is also the cheapest feasible mixture over all channels.
The second cheapest, \(\{\)portal, clinics South\(\}\), has margin \(0.030\)
and costs \(\$95.88\). The same support is cheapest for \(c_0=0.01\)
(\(45\) supports). For \(c_0=0.05\) and \(0.10\), only \(12\) and \(2\)
supports qualify, all offline at \(\$199\). The cheap supports have small
margins because the online and portal channels are close to the target
(\(13.9\%\) and \(13.4\%\) Black).

\paragraph{The myopic rule on this target.} To correct a deficit of Black
recruits, the myopic rule chooses clinics South (\(24.5\%\) Black), as
\(I^\star\) does. To correct an excess, online recruits move the share only
by \(0.016\) per recruit, while Spanish-language clinics (\(0.4\%\) Black)
move it by \(0.151\), which is more per dollar. The rule therefore settles
on \(\{\)clinics South, Spanish-language\(\}\) (\(60\%\) and \(36\%\) of
recruits at \(T=10^5\)), a support with the largest margin, \(0.128\), and
the highest cost, \(\$199\). This is the situation of
Example~\ref{ex:myopic}: a cheap source close to the target corrects slowly,
and the myopic rule replaces it by an expensive one.

\paragraph{Cost and error against the sample size.}
Figure~\ref{fig:app-budget}(a,b) and Table~\ref{tab:cost} report the cost
per recruit, i.e., the total cost divided by the sample size \(T\),
exploration included, and the error \(\|\widehat p_T-p_G\|_2\), averaged
over 200 runs. Algorithm~\ref{alg:cost-support}\(\to\)\ref{alg:ucb-sampling}
with \(n=1{,}600\) and \(c_0=0.02\) explores for \(25{,}600\) recruits at
\(\$183.13\) each (\(\$4.69\)M); its cost per recruit then falls to
\(\$123.26\) at \(T=10^5\) and \(\$105.51\) at \(T=3\times10^5\), towards
the cost of the selected support. Of the 200 runs, \(99\) select \(I^\star\),
\(95\) select \(\{\)portal, clinics South\(\}\), five select an offline pair,
and one selects \(\{\)Medicare, Spanish-language\(\}\), whose true hull
misses \(p_G\); that run stalls at an error of \(1.2\times10^{-2}\) and
accounts for most of the mean error (\(1.1\times10^{-4}\) at \(T=10^5\),
against a median of \(3.5\times10^{-5}\)). The myopic rule pays less early
on, while it samples the cheap channels to learn their means (\(\$149\) at
\(T\approx10^3\)), but more than
Algorithm~\ref{alg:cost-support}\(\to\)\ref{alg:ucb-sampling} from
\(T\approx10^4\) on, and \(\$197.08\) at \(T=10^5\). Its error equals that
of Algorithm~\ref{alg:ucb-sampling} on all channels and of the oracle closed
loop (\(1.1\times10^{-5}\) at \(T=10^5\)): with margin \(0.128\), its
imbalance \(\|S_T\|\) is about \(1\), against about \(6\) for the closed loop
on \(I^\star\) with known means (\(5.8\times10^{-5}\) at \(T=10^5\)). With
\(n=800\) and \(n=400\), exploration is cheaper, but \(4\%\) and \(12\%\) of
runs select a support that misses \(p_G\)
(Table~\ref{tab:selection}), and the mean error stalls at
\(3.7\times10^{-4}\) and \(2.5\times10^{-3}\).

\paragraph{Error at a fixed budget.} Figure~\ref{fig:app-budget}(c) reports
the error after a total budget between \(\$1\)M and \(\$16\)M. With
\(\$16\)M, Algorithm~\ref{alg:cost-support}\(\to\)\ref{alg:ucb-sampling}
(\(n=1{,}600\)) recruits \(140{,}100\) participants (\(\$114\) each) and
reaches an error of \(2.8\times10^{-5}\) (median \(1.6\times10^{-5}\)); the
myopic rule recruits \(81{,}300\) (\(\$197\) each) and reaches
\(1.2\times10^{-5}\) (median \(8.5\times10^{-6}\)), as does
Algorithm~\ref{alg:ucb-sampling} on all channels. Even with known means, the
closed loop on \(I^\star\) recruits \(175{,}500\) participants and reaches
\(2.2\times10^{-5}\): at a fixed budget, the larger sample of the cheaper
support does not make up for its smaller margin. Below \(\$5\)M,
Algorithm~\ref{alg:cost-support} is still exploring.

\begin{figure}[!ht]
\centering
\includegraphics[width=\textwidth]{figA_budget}
\caption{Cost-aware sampling for the Black/other target, online channels at
\(\$72\) and offline channels at \(\$199\) (mean over 200 runs). (a) Cost
per recruit, exploration included; dotted: \(B(I^\star)\). (b) Error
against the sample size. (c) Error at a fixed total budget.
Algorithm~\ref{alg:cost-support} uses \(c_0=0.02\) and \(n\) exploration
samples per channel; Algorithm~\ref{alg:ucb-sampling}, also after
Algorithm~\ref{alg:cost-support}, and the myopic rule use radius
\(0.3\,r_a\).}
\label{fig:app-budget}
\end{figure}

\begin{table}[!ht]
\caption{Cost-aware sampling for the Black/other target at \(T=10^5\) (200
replications). Algorithm~\ref{alg:ucb-sampling}, also after
Algorithm~\ref{alg:cost-support} (\(c_0=0.02\)), and the myopic rule use
radius \(0.3\,r_a\); Algorithm~\ref{alg:cost-support} never returned
\textsc{Fail}.}
\label{tab:cost}
\centering\small
\begin{tabular}{llrr}
\hline
policy & means & \(\|\widehat p_T-p_G\|_2\) & \(\$\)/recruit\\
\hline
online only & -- & \(2.3\cdot10^{-2}\) & 72.00\\
uniform & -- & \(2.1\cdot10^{-2}\) & 183.13\\
oracle open loop on the cheapest mixture & known & \(1.2\cdot10^{-3}\) & 91.27\\
oracle closed loop on \(I^\star\) & known & \(5.8\cdot10^{-5}\) & 91.32\\
oracle closed loop on all channels & known & \(1.1\cdot10^{-5}\) & 199.00\\
Alg.~\ref{alg:cost-support}\(\to\)\ref{alg:ucb-sampling}, \(n=400\) & learnt & \(2.6\cdot10^{-3}\) & 107.21\\
Alg.~\ref{alg:cost-support}\(\to\)\ref{alg:ucb-sampling}, \(n=800\) & learnt & \(3.8\cdot10^{-4}\) & 113.53\\
Alg.~\ref{alg:cost-support}\(\to\)\ref{alg:ucb-sampling}, \(n=1{,}600\) & learnt & \(1.1\cdot10^{-4}\) & 123.26\\
myopic cost-ratio rule & learnt & \(1.1\cdot10^{-5}\) & 197.08\\
Algorithm~\ref{alg:ucb-sampling} on all channels & learnt & \(1.1\cdot10^{-5}\) & 198.79\\
\hline
\end{tabular}
\end{table}

\paragraph{The myopic rule on Example~\ref{ex:myopic}.} A simulation of
Example~\ref{ex:myopic} (100 replications, \(T=10^5\)) confirms the
analysis. The myopic rule with the tie rule of
Algorithm~\ref{alg:sign-based} samples \(B\) and \(C\) in equal proportions
and pays \(2.00\) per recruit; the adaptive rule on \(\{A,B\}\) pays
\(1.00\); and Algorithm~\ref{alg:cost-support} (\(n=1{,}000\), \(c_0=0.1\))
followed by Algorithm~\ref{alg:ucb-sampling} pays \(1.02\), the excess
being the exploration of \(C\). Because \(p_G=0.5\), \(S_t\) often returns
to zero in this example, so the tie rule matters: breaking ties towards the
least-sampled source lowers the cost of the myopic rule to \(1.62\), still
\(62\%\) above the support-based rule. With estimated means (radius
\(0.3\,r_a\)), the myopic rule pays \(1.61\). On the NHIS instance \(S_t\)
rarely returns to zero, and the tie rule plays no role.

\paragraph{Support selection.} Table~\ref{tab:selection} and
Figure~\ref{fig:app-cost}(a,b) run Algorithm~\ref{alg:cost-support} alone
(250 replications per \(n\)). The main risk at small \(n\) is selecting a
support whose true hull does not contain \(p_G\): then
Algorithm~\ref{alg:ucb-sampling} cannot reach the target, and the error
stalls at the distance from \(p_G\) to that hull. This happens in
\(54\%\) of runs at \(n=25\) and \(c_0=0.02\), and in none at
\(n=1{,}600\); with costs drawn log-uniformly from the meta-analysis ranges
(last column), the rate is \(3\%\) at \(n=1{,}600\). The guarantee of
Theorem~\ref{thm:cost-support}, \(c(\widehat I)\ge c_0/2\), holds with
probability \(0.95\) at \(n=800\) and \(1.00\) at \(n=1{,}600\) for
\(c_0=0.02\), whereas the theorem asks for \(n\ge286{,}185\) (and
\(1{,}144{,}739\) and \(45{,}790\) at \(c_0=0.01\) and \(0.05\)). Its cost
bound at \(n=1{,}600\) is \(\$1{,}996\) per recruit, vacuous at these
budgets. Exact recovery of \(I^\star\) is rare (\(44\%\) of runs at
\(n=1{,}600\) and \(c_0=0.02\)), because \(\{\)portal, clinics South\(\}\)
costs only \(\$4.61\) more and has a larger margin; the selected supports
that contain \(p_G\) cost \(\$97.21\) on average, \(\$5.94\) above
\(B(I^\star)\). For \(c_0=0.05\), where \(B(I^\star)=\$199\), the selected
supports are cheaper than \(I^\star\) on average, as the theorem allows a
support with \(c_0/2\le c(I)<c_0\).

\begin{table}[!ht]
\caption{Support selection (Algorithm~\ref{alg:cost-support}, 250
replications per \(n\)): probability that the selected support does not
contain \(p_G\) in its true hull (``infeas.''), and probability that
\(c(\widehat I)\ge c_0/2\) (among runs with an admissible support). The
last column averages the infeasibility rate for \(c_0=0.02\) over the base
costs and nine cost vectors drawn log-uniformly from the meta-analysis
ranges.}
\label{tab:selection}
\centering\small
\begin{tabular}{r|rr|rr|rr|r}
\hline
 & \multicolumn{2}{c|}{\(c_0=0.01\)} & \multicolumn{2}{c|}{\(c_0=0.02\)} & \multicolumn{2}{c|}{\(c_0=0.05\)} & \(c_0=0.02\)\\
\(n\) & infeas. & \(c(\widehat I)\ge\frac{c_0}{2}\) & infeas. & \(c(\widehat I)\ge\frac{c_0}{2}\) & infeas. & \(c(\widehat I)\ge\frac{c_0}{2}\) & 10 cost vectors\\
\hline
25 & 0.54 & 0.44 & 0.54 & 0.42 & 0.54 & 0.14 & 0.47\\
50 & 0.46 & 0.52 & 0.46 & 0.52 & 0.33 & 0.28 & 0.42\\
100 & 0.39 & 0.60 & 0.39 & 0.59 & 0.32 & 0.28 & 0.34\\
200 & 0.32 & 0.68 & 0.26 & 0.73 & 0.16 & 0.42 & 0.25\\
400 & 0.18 & 0.82 & 0.12 & 0.86 & 0.06 & 0.57 & 0.13\\
800 & 0.06 & 0.94 & 0.04 & 0.95 & 0.00 & 0.74 & 0.06\\
1{,}600 & 0.01 & 0.99 & 0.00 & 1.00 & 0.00 & 0.85 & 0.03\\
\hline
\end{tabular}
\end{table}

\paragraph{The margin threshold \(c_0\).} Table~\ref{tab:c0} and
Figure~\ref{fig:app-cost}(c) sweep \(c_0\) and \(n\) at \(T=10^5\). A larger
\(c_0\) buys accuracy with cost: for \(n=800\), the cost rises from
\(\$107.6\) at \(c_0=0.01\) (error \(1.5\times10^{-3}\)) to \(\$197.0\) at
\(c_0=0.10\) (error \(1.3\times10^{-5}\)), because only offline supports
have margin at least \(0.05\). Algorithm~\ref{alg:cost-support} never
returned \textsc{Fail}. A longer exploration (\(n=1{,}600\), \(26\%\) of
the sample) lowers the error for \(c_0\le0.02\) by a factor of \(74\) to
\(103\) relative to \(n=400\), as it avoids infeasible supports, but raises
the cost, because exploration is spread uniformly over all channels.

\begin{table}[!ht]
\caption{Algorithm~\ref{alg:cost-support}\(\to\)\ref{alg:ucb-sampling}
(radius \(0.3\,r_a\)) for the Black/other target at \(T=10^5\) (100
replications): error and cost per recruit.}
\label{tab:c0}
\centering\small
\begin{tabular}{r|rr|rr|rr}
\hline
 & \multicolumn{2}{c|}{\(n=400\)} & \multicolumn{2}{c|}{\(n=800\)} & \multicolumn{2}{c}{\(n=1{,}600\)}\\
\(c_0\) & error & \(\$\) & error & \(\$\) & error & \(\$\)\\
\hline
0.01 & \(4.4\cdot10^{-3}\) & 99.8 & \(1.5\cdot10^{-3}\) & 107.6 & \(4.3\cdot10^{-5}\) & 120.8\\
0.02 & \(3.2\cdot10^{-3}\) & 106.9 & \(8.3\cdot10^{-4}\) & 113.8 & \(4.3\cdot10^{-5}\) & 124.1\\
0.05 & \(4.4\cdot10^{-4}\) & 145.4 & \(2.4\cdot10^{-4}\) & 161.3 & \(2.5\cdot10^{-5}\) & 172.2\\
0.10 & \(1.4\cdot10^{-5}\) & 194.1 & \(1.3\cdot10^{-5}\) & 197.0 & \(1.4\cdot10^{-5}\) & 194.9\\
\hline
\end{tabular}
\end{table}

\begin{figure}[!ht]
\centering
\includegraphics[width=\textwidth]{figA_cost}
\caption{Support selection and the margin threshold, Black/other target.
(a,b) Support selection as in Table~\ref{tab:selection}. (c) Error against
cost of Algorithm~\ref{alg:cost-support}\(\to\)\ref{alg:ucb-sampling} at
\(T=10^5\); \(c_0\in\{0.01,0.02,0.05,0.10\}\) increases from left to right
on every curve.}
\label{fig:app-cost}
\end{figure}

\FloatBarrier
\subsection{Downstream Estimates and Reweighting}
\label{app:exp-downstream}

Composition matching controls the groups in the target, not other
variables. In a respondent-level simulation for the ethnic target, each
recruit carries six health variables from NHIS. Table~\ref{tab:downstream}
reports the bias of their sample means after \(2\times10^4\) recruits
(standard deviations across replications are at most \(0.7\) pp).
Algorithm~\ref{alg:ucb-sampling} (radius \(0.3\,r_a\)) matches the ethnic
composition to within \(1.1\times10^{-4}\), so its remaining bias arises
within groups: the
channels it relies on over-represent, within each ethnic group, adults in
poorer health. For five of the six variables its absolute bias exceeds
that of online-only recruitment. Post-stratification on the same three
groups also corrects only the group proportions, and so it cannot remove
this within-group bias either. As a correction of the estimate,
reweighting is inexpensive here: post-stratifying an online-only sample
has design effect \(1.03\), i.e., it loses \(3\%\) of the effective sample
size. It cannot, however, enrol participants who were never recruited.
The target vector should therefore include the variables that an
endpoint depends on; by Section~\ref{sec:experiments}, this lowers the
margin and hence the accuracy achievable with a given set of channels.

\begin{table}[!ht]
\caption{Bias (pp) of sample means after \(2\times10^4\) recruits in the
respondent-level simulation (200 replications), and Kish design effect of
post-stratification on the three ethnic groups.}
\label{tab:downstream}
\centering\small
\setlength{\tabcolsep}{4.5pt}
\begin{tabular}{lrrrrrrr}
\hline
 & insulin & hypertension & heart disease & obesity & disability & fair/poor health & design effect\\
\hline
population value (\%) & 31.9 & 73.1 & 16.3 & 51.4 & 23.8 & 42.8 & \\
Alg.~\ref{alg:ucb-sampling}, radius \(0.3\,r_a\) & \(+1.7\) & \(+2.4\) & \(+1.5\) & \(+2.7\) & \(+3.2\) & \(+4.1\) & 1.000\\
oracle closed loop & \(+1.0\) & \(+1.7\) & \(+0.6\) & \(+4.1\) & \(+2.1\) & \(+2.6\) & 1.000\\
online only & \(-0.1\) & \(-0.9\) & \(-0.5\) & \(+4.4\) & \(-2.9\) & \(-3.6\) & 1.030\\
uniform & \(+2.5\) & \(+2.1\) & \(+2.6\) & \(-0.1\) & \(+3.7\) & \(+5.5\) & 1.006\\
\hline
\end{tabular}
\end{table}

\FloatBarrier

\end{document}